\documentclass[11pt, a4paper]{report}
% IEEE body text = 10pt. Using report class for chapter support.

% =============================================================================
% PACKAGES
% =============================================================================

% --- Encoding & Language ---
\usepackage[utf8]{inputenc}
\usepackage[T1]{fontenc}
\usepackage[british]{babel}

% --- Typography ---
% IEEE uses Times Roman throughout
\usepackage{mathptmx}              % Times Roman for text + math
\usepackage[scaled=0.90]{helvet}   % Helvetica (sans-serif fallback)
\usepackage{courier}               % Courier for monospace / code

% --- Page Geometry (IEEE A4 spec) ---
% IEEE A4: side = 13mm, top = 19mm, bottom = 43mm
% Single-column width = 184mm
\usepackage[
    a4paper,
    left=13mm,
    right=13mm,
    top=19mm,
    bottom=43mm,
    headheight=12pt,
    headsep=10pt,
    footskip=20pt
]{geometry}

% --- Line Spacing ---
% IEEE body text: 10pt font / 12pt line spacing = factor of 1.2
% LaTeX default for 10pt is ~12pt baselineskip, so single spacing is correct.
\usepackage{setspace}
\singlespacing

\usepackage[parfill]{parskip}

% --- Headers & Footers ---
\usepackage{fancyhdr}
\pagestyle{fancy}
\fancyhf{}
\fancyhead[L]{\footnotesize\nouppercase{\leftmark}}
\fancyhead[R]{\footnotesize\thepage}
\renewcommand{\headrulewidth}{0.4pt}
\fancypagestyle{plain}{%
    \fancyhf{}
    \fancyfoot[C]{\footnotesize\thepage}
    \renewcommand{\headrulewidth}{0pt}
}

% --- Graphics & Figures ---
\usepackage{graphicx}
\usepackage{float}
\usepackage{subcaption}
\usepackage{chngcntr}

% IEEE Figure captions: 8pt, 9pt line spacing, centred below figure
% IEEE Table captions: 8pt, 9pt line spacing, centred above table, small caps label
\usepackage[
    font={footnotesize, singlespacing},  % 8pt at 10pt base
    labelfont=bf,
    labelsep=period,
    justification=centering,
    tableposition=top,
    figureposition=bottom
]{caption}

% Override table caption label to use small caps "TABLE" with Arabic numerals
\DeclareCaptionLabelFormat{ieeetable}{\textsc{Table~#2}}
\captionsetup[table]{labelformat=ieeetable, labelsep=newline, justification=centering}

% Figure numbering: Arabic, label "Fig."
\DeclareCaptionLabelFormat{ieeefigure}{Fig.~#2}
\captionsetup[figure]{labelformat=ieeefigure, labelsep=period, justification=centering}

% Tie Table and Figure numbering to Chapters (e.g., Table 1.1, Fig 2.1)
\counterwithin{table}{chapter}
\counterwithin{figure}{chapter}

% --- Tables ---
\usepackage{booktabs}
\usepackage{array}
\usepackage{multirow}
\usepackage{tabularx}

% --- Mathematics ---
\usepackage{amsmath}
\usepackage{amssymb}
\usepackage{bm}
\usepackage{braket}

% --- Colours ---
\usepackage{xcolor}
\definecolor{codeBackground}{RGB}{248,248,250}
\definecolor{codeFrame}{RGB}{190,190,200}
\definecolor{codeKeyword}{RGB}{0,80,160}
\definecolor{codeComment}{RGB}{80,128,80}
\definecolor{codeString}{RGB}{160,48,48}
\definecolor{codeNumber}{RGB}{140,140,150}
\definecolor{linkBlue}{RGB}{0,60,120}

% --- Code Listings ---
\usepackage{listings}
\lstdefinestyle{ieeeCode}{
    backgroundcolor=\color{codeBackground},
    frame=single,
    rulecolor=\color{codeFrame},
    framesep=3pt,
    xleftmargin=6pt,
    xrightmargin=4pt,
    % IEEE code should be small --- 8pt matches footnote/caption size
    basicstyle=\ttfamily\scriptsize\linespread{1.0}\selectfont,
    keywordstyle=\color{codeKeyword}\bfseries,
    commentstyle=\color{codeComment}\itshape,
    stringstyle=\color{codeString},
    numberstyle=\tiny\color{codeNumber}\ttfamily,
    numbers=left,
    numbersep=6pt,
    breaklines=true,
    breakatwhitespace=false,
    showstringspaces=false,
    captionpos=b,
    tabsize=4,
    columns=flexible,
    keepspaces=true,
    aboveskip=10pt,
    belowskip=6pt,
    abovecaptionskip=4pt,
    belowcaptionskip=2pt,
    literate={-}{-}1
}
\lstset{style=ieeeCode}
\renewcommand{\lstlistingname}{Listing}
\renewcommand{\lstlistlistingname}{List of Listings}

% --- Algorithms ---
\usepackage[ruled, vlined, linesnumbered]{algorithm2e}
\SetAlgoLined
\SetKwComment{Comment}{$\triangleright$\ }{}
\SetAlCapSkip{6pt}
% Match algorithm caption size to IEEE caption spec (8pt)
\SetAlCapFnt{\footnotesize}
\SetAlCapNameFnt{\footnotesize\bfseries}

% --- Tie Counters to Chapters ---
% MUST BE PLACED HERE (After the packages are loaded)
\counterwithin{algocf}{chapter}     % This is the counter used by algorithm2e
\counterwithin{equation}{chapter}   % Optional but highly recommended!

% The listings package delays creating its counter until \begin{document}
\AtBeginDocument{
    \counterwithin{lstlisting}{chapter}
}

% --- Hyperlinks ---
\usepackage[
    colorlinks=true,
    linkcolor=black,
    citecolor=linkBlue,
    urlcolor=linkBlue,
    bookmarks=true,
    bookmarksnumbered=true,
    pdfauthor={Om Kalyan},
    pdftitle={Beyond Classical Cryptography: Simulating and Benchmarking QKD},
    pdfsubject={BEng Dissertation},
    pdfkeywords={QKD, BB84, B92, E91, Qiskit, Post-Quantum Cryptography}
]{hyperref}
\usepackage[noabbrev, capitalise]{cleveref}

% --- Bibliography ---
% IEEE references: 8pt, 9pt line spacing, bracketed Arabic numerals
\usepackage{cite}

% --- Glossary, Acronyms & Nomenclature ---
% Requires a build step: run `makeglossaries <jobname>` between pdflatex passes
\usepackage[
    acronym,         % creates a separate acronym list (e.g. QKD, BB84)
    symbols,         % creates a separate symbols/nomenclature list
    toc,             % adds each printed glossary to the Table of Contents
    nogroupskip,     % suppresses the letter-group gaps (A, B, C) for a cleaner look
    nopostdot,       % removes the trailing full stop after each description
    nonumberlist     % suppresses back-references to page numbers
]{glossaries-extra}

% Define the custom symbols database BEFORE makeglossaries
%\newglossary[slg]{symbols}{sym}{sbl}{Nomenclature}

\loadglsentries{glossary.tex}

\makeglossaries

% --- Lists ---
\usepackage{enumitem}
\setlist{nosep, leftmargin=*}

% --- Paragraph formatting ---
% IEEE uses indented paragraphs with no extra inter-paragraph space
\setlength{\parindent}{1em}
\setlength{\parskip}{2pt}

% =============================================================================
% HEADING STYLES --- IEEE Conference Manuscript Format Definition
% =============================================================================
\usepackage{titlesec}

% --- CHAPTER ─────────────────────────────────────────────────────────────────
% Dissertation addition (IEEE conf. papers have no chapters).
% Styled as a prominent centred heading, consistent with IEEE aesthetic.
% Uses 14pt bold small caps, centred, with "CHAPTER I" label above.
\titleformat{\chapter}[display]
    {\normalfont\fontsize{15}{17}\bfseries\scshape\centering}
    {\MakeUppercase{\chaptertitlename}~\thechapter}
    {7pt}
    {\fontsize{15}{17}\selectfont\bfseries}
\titlespacing*{\chapter}{0pt}{-10pt}{20pt}

% --- SECTION ─────────────────────────────────────────────────────────────────
\renewcommand{\thesection}{\thechapter.\arabic{section}} % e.g., 1.1, 1.2
\titleformat{\section}
    {\normalfont\fontsize{11}{15}\selectfont\scshape\centering}
    {\thesection}{1em}{} % Removed the trailing dot after the number

% --- SUBSECTION ──────────────────────────────────────────────────────────────
\renewcommand{\thesubsection}{\thesection.\arabic{subsection}} % e.g., 1.1.1
\titleformat{\subsection}
    {\normalfont\fontsize{11}{13}\selectfont\itshape}
    {\thesubsection}{1em}{} % Removed the trailing dot

% --- SUBSUBSECTION ───────────────────────────────────────────────────────────
\renewcommand{\thesubsubsection}{\thesubsection.\arabic{subsubsection}} % e.g., 1.1.1.1
\titleformat{\subsubsection}[runin]
    {\normalfont\fontsize{11}{13}\selectfont\itshape}
    {\thesubsubsection}{1em}{}[:] % Removed the trailing parenthesis

% --- PARAGRAPH (IEEE "Heading 4") ───────────────────────────────────────────
% Spec: 10pt italic, left-justified,
%       Lowercase alpha + right paren + colon, run-in, indent 2em
\renewcommand{\theparagraph}{\alph{paragraph}}
\titleformat{\paragraph}[runin]
    {\normalfont\fontsize{11}{13}\selectfont\itshape}
    {\theparagraph)}{0.5em}{}[:]
\titlespacing*{\paragraph}{2em}{6pt}{0.5em}

% Enable numbering depth through paragraph level
\setcounter{secnumdepth}{4}
\setcounter{tocdepth}{3}

% =============================================================================
% CUSTOM COMMANDS (QKD Notation)
% =============================================================================

\newcommand{\QBER}{\ensuremath{\mathrm{QBER}}}
\newcommand{\keyrate}{\ensuremath{R_{\mathrm{sec}}}}

% =============================================================================
% FOOTNOTE SIZE (IEEE: 8pt, 9pt line spacing)
% =============================================================================
\renewcommand{\footnotesize}{\fontsize{9}{10}\selectfont}

% glossary addition

%% =============================================================================
% GLOSSARY ENTRIES
% =============================================================================

% --- ACRONYMS ----------------------------------------------------------------
% Usage in text: \gls{qkd}  first use: "Quantum Key Distribution (QKD)"
% subsequent: "QKD"
% \glspl{qkd}  plural form

\newacronym{qkd}{QKD}{Quantum Key Distribution}
\newacronym{bb84}{BB84}{Bennett--Brassard 1984 Protocol}
\newacronym{b92}{B92}{Bennett 1992 Protocol}
\newacronym{e91}{E91}{Ekert 1991 Protocol}
\newacronym{qber}{QBER}{Quantum Bit Error Rate}
\newacronym{pqc}{PQC}{Post-Quantum Cryptography}
\newacronym{otp}{OTP}{One-Time Pad}
\newacronym{prng}{PRNG}{Pseudo-Random Number Generator}
\newacronym{epe}{EPE}{Entanglement Purification via Eavesdropping}
\newacronym{des}{DES}{Data Encryption Standard}
\newacronym{aes}{AES}{Advanced Encryption Standard}
\newacronym{chsh}{CHSH}{Clauser--Horne--Shimony--Holt}
\newacronym{edi}{EDI}{Equality, Diversity, and Inclusion}
\newacronym{epr}{EPR}{Einstein--Podolsky--Rosen}
\newacronym{cv}{CV}{Continuous-Variable}
\newacronym{dv}{DV}{Discrete-Variable}
\newacronym[sort=pandm]{pandm}{P\&M}{Prepare-and-Measure}
\newacronym{eb}{EB}{Entanglement-Based}
\newacronym{mdi}{MDI}{Measurement-Device-Independent}
\newacronym{cc}{CC}{Classical Channel}
\newacronym{oop}{OOP}{Object-Oriented Programming}
\newacronym{bsc}{BSC}{Binary Symmetric Channel}
\newacronym{sdk}{SDK}{Software Development Kit}
\newacronym{ide}{IDE}{Integrated Development Environment}
\newacronym{qchan}{QC}{Quantum Channel}
\newacronym{qcomp}{QC}{Quantum Computer}
\newacronym{skr}{SKR}{Secure Key Rate}
\newacronym{p2p}{P2P}{Peer-to-Peer}

% --- SYMBOLS (Nomenclature) --------------------------------------------------
% Usage in text: \gls{sym:qber}  renders the symbol with a hyperlink
%                \glsdesc{sym:qber} renders just the description

\newglossaryentry{sym:qber}{
    type    = symbols,
    name    = {\ensuremath{\mathrm{QBER}}},
    description = {Quantum Bit Error Rate --- ratio of erroneous bits to total
                   bits received over the quantum channel},
    sort    = {QBER}
}

\newglossaryentry{sym:rsec}{
    type    = symbols,
    name    = {\ensuremath{R_{\mathrm{sec}}}},
    description = {Secret-key generation rate (bits per channel use)},
    sort    = {Rsec}
}

\newglossaryentry{sym:theta}{
    type    = symbols,
    name    = {\ensuremath{\theta}},
    description = {Polarisation angle of a photon (radians)},
    sort    = {theta}
}

\newglossaryentry{sym:epsilon}{
    type    = symbols,
    name    = {\ensuremath{\varepsilon}},
    description = {Security parameter --- maximum acceptable eavesdropper
                   information leakage},
    sort    = {epsilon}
}

% --- GENERAL GLOSSARY TERMS --------------------------------------------------
% Usage in text: \gls{entanglement}

\newglossaryentry{entanglement}{
    name        = {entanglement},
    description = {A quantum mechanical phenomenon in which two or more particles
                   share a correlated quantum state, such that measuring one
                   particle instantaneously determines properties of the other
                   regardless of the distance separating them}
}

\newglossaryentry{qubit}{
    name        = {qubit},
    description = {The fundamental unit of quantum information; a two-level
                   quantum system that can exist in a superposition of states
                   \ensuremath{\ket{0}} and \ensuremath{\ket{1}}}
}

\newglossaryentry{sifting}{
    name        = {sifting},
    description = {The post-transmission step in QKD in which Alice and Bob
                   publicly compare measurement bases and discard bits where
                   incompatible bases were used}
}

\newglossaryentry{reconciliation}{
    name        = {reconciliation},
    description = {An error-correction stage in QKD that resolves discrepancies
                   between Alice's and Bob's sifted keys over an authenticated
                   classical channel}
}

\newglossaryentry{amplification}{
    name        = {privacy amplification},
    description = {A post-processing step that compresses the reconciled key to
                   reduce any partial information an eavesdropper may have
                   obtained, producing a shorter but provably secure final key}
}


% =============================================================================
% ENTRIES ADDED FOR: Introduction paragraph
% =============================================================================

% --- New Acronyms ------------------------------------------------------------

\newacronym{rsa}{RSA}{Rivest--Shamir--Adleman}
% First use in text: "Rivest--Shamir--Adleman (RSA)"
% Subsequent uses:   "RSA"

\newacronym{ecc}{ECC}{Elliptic Curve Cryptography}
% First use in text: "Elliptic Curve Cryptography (ECC)"
% Subsequent uses:   "ECC"

% --- New General Terms -------------------------------------------------------

\newglossaryentry{intfact}{
    name        = {integer factorisation},
    description = {The computational problem of decomposing a composite integer
                   into its prime factors. No classical algorithm is known to
                   solve this in polynomial time for large integers, forming the
                   security basis of RSA}
}

\newglossaryentry{disclog}{
    name        = {discrete logarithm},
    description = {The problem of finding the exponent \ensuremath{x} such that
                   \ensuremath{g^x \equiv h \pmod{p}} for given \ensuremath{g},
                   \ensuremath{h}, and prime \ensuremath{p}. Believed to be
                   classically hard; underlies the security of ECC and
                   Diffie--Hellman key exchange}
}

\newglossaryentry{shors}{
    name        = {Shor's algorithm},
    description = {A quantum algorithm, published by Peter Shor in 1994, that
                   solves integer factorisation and the discrete logarithm in
                   polynomial time on a fault-tolerant quantum computer,
                   breaking RSA and ECC}
}

\newglossaryentry{itsec}{
    name        = {information-theoretic security},
    description = {A security guarantee that holds against an adversary with
                   unlimited computational power, relying on mathematical
                   impossibility rather than assumed hardness of a problem.
                   QKD achieves this via the laws of quantum mechanics}
}

\newglossaryentry{qea}{
    name        = {Quantum Economic Advantage},
    description = {The threshold at which the practical cost of executing a
                   quantum attack --- accounting for hardware, error correction,
                   and runtime --- becomes economically feasible for a
                   sufficiently motivated adversary}
}

   % load all entry definitions


% counterwithin updates fro algirthims, listings, etc..


% =============================================================================
% DOCUMENT
% =============================================================================
\begin{document}

% ─────────────────────────────────────────────────────────────────────────────
% TITLE PAGE
% ─────────────────────────────────────────────────────────────────────────────
\begin{titlepage}
    \centering
    \thispagestyle{empty}
    \vspace*{3cm}

    % --- Title: IEEE spec = 24pt regular, 28pt line spacing, centred ---
    {\fontsize{28}{32}\selectfont\bfseries
    Quantum Key Distribution:\par
    A Simulation Analysis of\par
    BB84, B92 and E91\par}

    %Quantum Key Distribution: Simulating and Benchmarking BB84, B92, and E91

    \vspace{1.0cm}

        % --- Author Affiliation: IEEE spec = 10pt regular, 11pt line spacing, centred ---
    {\fontsize{14}{15}\selectfont
    A dissertation submitted for the degree of\par}

    \vspace{0.5cm}

    {\fontsize{14}{15}\selectfont\bfseries
    Bachelor of Engineering\\
    Electronics \& Computer Engineering\par}

    \vspace{1.0cm}

    % --- Author Name: IEEE spec = 11pt regular, 12pt line spacing, centred ---
    {\fontsize{14}{15}\selectfont
    Om Kalyan\par}

    \vspace{0.5cm}

    {\fontsize{10}{11}\selectfont
    230733440\par}

    \vspace{1.0cm}

    % --- University Crest ---
    
\includegraphics[width=0.24\textwidth]{misc images/ncl logo vector.png}

    \vspace{1.0cm}

    % --- Institution: 11pt to match author name tier ---
    {\fontsize{11}{13}\selectfont
    School of Engineering\\
    Newcastle University\par}

    \vspace{0.6cm}

    % --- Date ---
    {\fontsize{10}{11}\selectfont May 2026\par}

    \vspace{1cm}
    
    % word count
    {\fontsize{10}{11}\selectfont Word Count: 6668\par}


\end{titlepage}

% -------------------        cover pages      -------------------
% ─────────────────────────────────────────────────────────────────────────────
% PAGE 2: SELF-REFLECTION AND FEEDBACK
% ─────────────────────────────────────────────────────────────────────────────
\clearpage
\thispagestyle{empty}

\noindent\textbf{Self-reflection and feedback}

\vspace{1em}
\begin{itemize}
    \item \textbf{I believe a strong part of my submission is:}
    
The successful implementation and verification of three QKD protocols against existing literature. A key strength is the classical post-processing architecture used to translate raw quantum mechanics into a practical key distribution scheme. The highlight is the multivariable parameter sweeps, producing novel protocol performance plots not found in my literature survey.
    
    \vspace{3em}

    \item \textbf{An area I struggled with is:}
    
The initial steep learning curve of quantum mechanics, specifically the abstract themes and dense mathematical theory. I overcame this by methodically breaking down the theoretical physics into functional requirements, which allowed me to translate the abstract mathematical concepts into Python logic for the simulator.

    \vspace{3em}

    \item \textbf{I would particularly like feedback/help with:}
    
How to refine my presentation skills to effectively communicate dense, highly technical work. I would appreciate any advice on delivering the breadth of this project, particularly with balancing the complex mathematics and software engineering decisions in an accessible format. Going forward, any of your thoughts will be very valuable.

    \vspace{3em}
\end{itemize}

% ─────────────────────────────────────────────────────────────────────────────
% PAGE 3: ENGINEERING SKILLS PORTFOLIO INDEXING
% ─────────────────────────────────────────────────────────────────────────────
\clearpage
\thispagestyle{empty}

\noindent\textbf{Engineering skills portfolio indexing}

\vspace{1em}
\noindent\begin{tabular}{|p{0.72\textwidth}|p{0.20\textwidth}|}
\hline
\textbf{Engineering skill} & \textbf{Page(s) where demonstrated} \\
\hline
\textbf{Environmental and societal impact} Where in your report have you demonstrated that you have considered the wider impact on society to the broader area that you have been working in? For example, if your project linked to reduction of greenhouse gasses or personal mobility, how does your work fit into these broader societal goals? & 1, 4, 40. \\
\hline
\textbf{Ethical considerations} All projects contain ethical considerations, although some are more obvious than others. Have you researched any potential negative impacts of your project? Are there ethical questions in literature or journalism about the area you are working in? What is the general consensus of the engineering community of these issues? & 1, 2, 40\\
\hline
\textbf{Risk management and mitigation} For example, have you planned for unexpected outcomes or circumstances? How did you manage changes to your project? Did you decide to change your project and if so how did it change the risk profile of the work? This is usually achieved through regular review, reflection and modification of your project plan, but you must demonstrate this. & 8, 24, 39.\\
\hline
\textbf{Intellectual property security and data management} Whilst some sectors have a more obvious security aspect, all projects require a degree of data protection and management. How have you ensured that the data or intellectual property (IP) which you have generated within your project is secure from external actors or accidental loss? Have you used backups or password protected file stores? Where have you demonstrated an awareness of how IP is captured. How was information management included in part of your project planning? & 8, 9.\\
\hline
\textbf{Equality, diversity and inclusion} Linking to your considerations of societal impact, where have you considered the inclusivity of your work? Have you demonstrated that you have considered how your work may impact different groups and made efforts to minimise this? Is your work accessible to a wide range of users and could others from a different background to yourself repeat your results based on what you have written in your report? & 2, 40, \\
\hline
\textbf{Quality of information} Where in your report have you discussed the \textit{quality} of the information you have researched or the data you have presented? Have you reflected on how this could be improved? & 24, 25, 39. \\
\hline
\textbf{Continual professional development and self-reflection} Where in your report have you reflected on your own performance throughout the project? Have you identified personal strengths and achievements as well as weaknesses? Have you made plans for exploiting the strengths and improving the weaknesses? & 39, 40. \\
\hline
\end{tabular}



% ─────────────────────────────────────────────────────────────────────────────
% FRONT MATTER
% ─────────────────────────────────────────────────────────────────────────────
\pagenumbering{roman}

% --- Abstract ---
\clearpage
\chapter*{Abstract}
%\addcontentsline{toc}{chapter}{Abstract}

% IEEE Abstract: 9pt bold, "Abstract---" bold italic, 10pt line spacing
{\fontsize{9}{10}\selectfont\bfseries

This project develops a Python and Qiskit simulation platform to benchmark the \gls{bb84}, \gls{b92}, and \gls{e91} \gls{qkd} protocols under realistic channel conditions. A modular \gls{oop} architecture models state preparation, transmission, and post-processing under depolarising noise and intercept-resend attacks, with performance quantified via \gls{qber}, Mutual Information, and Devetak-Winter \gls{skr}. Outputs are validated against theoretical security bounds and published simulation benchmarks.

\ 

Multivariable analysis confirms \gls{bb84} maintains secure key generation up to its established $11\%$ \gls{qber} threshold, whilst \gls{b92}'s two-state encoding collapses at a stricter $6.5\%$ limit. Critically, despite sharing near-identical \gls{qber} trajectories under noise, \gls{e91} sustains a superior \gls{skr} above $p \approx 2.5\%$, demonstrating that \gls{eb} correlations provide a more favourable Devetak-Winter security penalty than state-based encoding under equivalent channel conditions. The \gls{e91} implementation further identifies two distinct noise thresholds: the Devetak-Winter bound forces \gls{skr} to zero at $P \approx 0.130$, preceding full \gls{chsh} classical bound failure at $P \approx 0.165$. Ultimately, the platform establishes the foundations for a digital-twin systems engineering approach to \gls{qkd} deployment evaluation.


\par
}

% --- Acknowledgements ---
\clearpage
\phantomsection
\chapter*{Acknowledgements}
\addcontentsline{toc}{chapter}{Acknowledgements}

I would like to thank my supervisors Martin Johnston and Oktay \c{C}etinkaya for their guidance over the course of this project.

% --- Table of Contents ---
\clearpage
\tableofcontents


% --- List of Figures ---
\clearpage
\phantomsection
\listoffigures
\addcontentsline{toc}{chapter}{List of Figures}

% --- List of Tables ---
\clearpage
\phantomsection
\listoftables
\addcontentsline{toc}{chapter}{List of Tables}

% --- List of Listings ---
\clearpage
\phantomsection
\lstlistoflistings
\addcontentsline{toc}{chapter}{List of Listings}

% --- List of Algorithms ---
\clearpage
\phantomsection
\listofalgorithms
\addcontentsline{toc}{chapter}{List of Algorithms}

% --- Acronyms ---
\printglossary[
    type  = acronym,
    title = {List of Acronyms},
    style = long              % two-column tabular: abbreviation | definition
]
\clearpage


% --- Nomenclature / Symbols ---
\glsaddall[types={symbols}] % Forces LaTeX to write the .sym file
\printglossary[
    type  = symbols,
    title = {Nomenclature},
    style = long
]
\clearpage
% --- General Glossary ---
\printglossary[
    title = {Glossary},
    style = altlist            % term on its own line, indented description below
]
\clearpage

%================================================================================
% MAIN CONTENT
%================================================================================
\clearpage
\pagenumbering{arabic}

% glossary reset
\glsresetall

%================================================================================
%                  Intro
%================================================================================
\chapter{Introduction}
%================================================================================

\section{A Brief History of Cryptology}

Since the dawn of days, the fundamental objective of cryptography has been to secure communication in the presence of adversaries. While early cryptology relied on manual ciphers and electromechanical systems \cite{Kahn1967Code, Hodges2014}, the discipline was formalised by Claude Shannon in 1945. Shannon's work on information theory and secrecy systems \cite{Shannon1945, Shannon1948} shifted cryptographic design away from physical obfuscation and towards rigorous mathematics. Crucially, Shannon established that perfect \gls{itsec} requires the \gls{otp}, which remains notoriously difficult to implement due to the challenge of secure key distribution.

As a practical compromise, modern digital security was built upon computational hardness rather than perfect secrecy. Standardised algorithms such as the \gls{des} (1977), the \gls{rsa} public-key cryptosystem (1978), and the \gls{aes} (2001) were developed to secure global networks \cite{DiffieHellman1976, RSA1978}. Today, these protocols underpin virtually all digital infrastructure, deriving security from the presumed difficulty of mathematical problems such as \gls{intfact} and the \gls{disclog}.


\section{The Quantum Threat to Classical Cryptography}

This reliance on computational hardness is now the fundamental vulnerability of classical public-key cryptography. \Gls{shors} \cite{Shor1997} demonstrated that a sufficiently powerful \gls{qcomp} could solve both integer factorisation and the discrete logarithm in polynomial time, rendering schemes such as \gls{rsa} and \gls{ecc} obsolete. As illustrated in \Cref{fig:quantumAdvantage}, estimates suggest that we are approaching \gls{qea}, where the resource cost of quantum attacks becomes economically feasible for motivated adversaries.

\begin{figure}[htbp]
    \centering
    \includegraphics[width=0.7\linewidth]{images.introduction/quantumAdvantage.png}
    \caption{The Quantum Economic Advantage.}
    \label{fig:quantumAdvantage}
\end{figure}

The ``harvest now, decrypt later'' attack model gives this threat immediate practical relevance. Adversaries can intercept classically encrypted data today and store it until a capable \gls{qcomp} becomes available \cite{NIST}. Consequently, data with long-term confidentiality requirements, such as state secrets and financial instruments, are already vulnerable.

To secure communications against quantum adversaries, two responses have emerged: \gls{pqc}, which operates entirely within classical hardware, and \gls{qkd}, which achieves true \gls{itsec}. \Gls{qkd} solves the historical key distribution problem of the \gls{otp} by exploiting quantum mechanical properties to distribute secure keys at scale, without relying on trusted couriers \cite{QKDbible}. This project focuses on the latter, developing a simulation platform to benchmark the \gls{bb84}, \gls{b92}, and \gls{e91} \gls{qkd} protocols.

\newpage
\section{Aims \& Objectives}

The primary aim of this project is to develop and evaluate a modular simulation platform for benchmarking \gls{qkd} protocols under realistic channel noise and active eavesdropping constraints. 

To achieve this aim, the project will pursue five core objectives:
\begin{enumerate}[label=\textbf{O\arabic*}, leftmargin=1.5cm]
    \item \textbf{Theoretical Foundation:} Synthesise the mechanics, security proofs, and evaluation metrics (\gls{qber}, \gls{skr}, Mutual Information) of the \gls{bb84}, \gls{b92}, and \gls{e91} protocols.
    \item \textbf{Platform Implementation:} Construct a Python/Qiskit simulator modelling state preparation, transmission, measurement, and classical post-processing.
    \item \textbf{Interference Modelling:} Implement tunable depolarising channel noise and active intercept-resend eavesdropping attacks to quantify protocol degradation.
    \item \textbf{Validation \& Benchmarking:} Extract secure key rates and performance thresholds through multivariable analysis, validating outputs against established theoretical bounds and literature precedents.
    \item \textbf{Broader Impact:} Contextualise \gls{qkd} deployment within wider ethical and \gls{edi} frameworks, evaluating the societal accessibility of quantum-safe infrastructure.
\end{enumerate}



%================================================================================
%                  Lit Review
%================================================================================
\chapter{Literature Review}
%================================================================================

\section{Early Work}
Presented in 1984, Bennett \& Brassard \cite{BB84} demonstrated that secure key exchange can be achieved through the use of non-orthogonal quantum states and basis incompatibility. This experiment proved that any eavesdropping attempt introduces detectable disturbances, laying the foundation for most of the research in the field of \gls{qkd}. The work of Ekert \cite{E91} later introduced the concept of leveraging the mechanics of \gls{entanglement} to verify security through Bell inequality violations. Bennett \cite{B92} continued to innovate in the field by simplifying the \gls{bb84} protocol to use only two non-orthogonal states through \gls{b92}, at the cost of efficiency. Taken together, these protocols represent three important design directions in \gls{qkd}: \gls{pandm} operation with strong foundational security arguments, \gls{eb} security verification, and reduced-state implementations that trade efficiency for simplicity.

\section{Protocol Classification}
% QKD whitepapers
\begin{figure}[htbp]
    \centering
    \includegraphics[width=0.6\linewidth]{images.lit/QKDclassification.png}
    \caption{Classification of QKD Protocol Schemes.}
    \label{fig:QKDclassification}
\end{figure}

Due to rapid developments in \gls{qkd}, protocols are commonly classified according to the encoding variable and the security architecture employed. In this context, \gls{cv} protocols encode information in quadrature observables such as amplitude and phase, whereas \gls{dv} protocols use finite-dimensional quantum states such as polarisation or \gls{entanglement} correlations, as depicted in \Cref{fig:QKDclassification}. \gls{cv} schemes therefore rely on continuous degrees of freedom of the quantum system and can often be implemented using standard optical hardware, although their practical noise tolerance depends on the channel and detection conditions. In comparison, \gls{dv} schemes use discrete quantum states, such as single photons or entangled photon pairs, as described by Kumar \textit{et al.} \cite{Kumar2025}. Protocols are also commonly grouped by their implementation architecture, typically \gls{pandm}, \gls{eb}, or \gls{mdi}, reflecting different trade-offs in implementation complexity and trust assumptions.


\newpage
\section{Experimental Studies}
Following the introduction of these protocols, a significant effort towards evaluating their performance under realistic channel conditions and imperfect devices has driven the need to understand detector inefficiencies, source imperfections, and finite-key effects prior to practical deployment \cite{Scarani2009, GLLP2004}. These studies bridged the gap between protocol design and implementable systems, motivating experimental demonstrations.

A major step towards practical security proof was the work of Shor and Preskill \cite{ShorPreskill2000}, who provided a simpler proof of unconditional security for \gls{bb84}. Their results strengthened confidence in \gls{qkd} as a viable secure communication method, encouraging further investigation into how protocol security degrades under realistic conditions.

\begin{figure}[htbp]
    \centering
    \includegraphics[width=0.7\linewidth]{images.lit/toshiba architecture for KMS.png}
    \caption{Toshiba Quantum Key Management System Topology.}
    \label{fig:KMS}
\end{figure}

Experimental work has also demonstrated that \gls{qkd} can be deployed over meaningful distances in optical fibre networks. For example, Boaron \textit{et al.} \cite{Boaron2018} reported secure \gls{qkd} over 421 km of optical fibre. These demonstrations not only validate the theoretical developments in the field, but also expose the practical constraints that must be accounted for before widespread adoption. Zhang's recent survey \cite{Zhang2025} indicates that the field is now moving towards larger-scale, integrated \gls{qkd} infrastructures to enable networked deployment and global connectivity. 

Commercial deployment by Toshiba Quantum across metropolitan and wider fibre networks (seen in \Cref{fig:KMS}) demonstrates that \gls{qkd} is no longer purely theoretical, but the practical use of these systems remains constrained by infrastructure compatibility, integration cost and operational complexity \cite{toshiba_qkd_dske, toshiba2023quantumsafe}. Despite this, the potential remains significant, with the deployment cases in \Cref{fig:applications} serving as viable use cases. 

\begin{figure}[htbp]
    \centering
    \includegraphics[width=0.7\linewidth]{images.lit/QKDapplicationsTNR.png}
    \caption{Real-World Applications of QKD.}
    \label{fig:applications}
\end{figure}

\section{Simulation Studies}
Simulation has become an important tool in \gls{qkd} research because it allows protocol behaviour to be evaluated before physical implementation. Rasay \textit{et al.} \cite{rasay2026} compare \gls{bb84}, \gls{b92}, and \gls{e91} under ideal and noisy conditions using \gls{qber}, showing that \gls{bb84} is generally more resilient than \gls{b92} when measurement uncertainty and channel noise are introduced. Devi \textit{et al.} \cite{devi2025} extend this idea by adding interactive visualisation of \gls{bb84} and \gls{e91}, but their work is still primarily comparative and educational rather than a unified benchmarking framework.

These studies collectively show that simulation is useful not only for demonstrating protocol principles, but also for revealing how protocol performance changes under common conditions. However, the literature still lacks a consistent framework in which \gls{bb84}, \gls{b92}, and \gls{e91} are evaluated using the same assumptions, noise model, and performance metrics.

Simulation toolkits such as Qiskit \cite{QiskitAer} and PennyLane \cite{pennylane2026docs} support quantum circuit modelling, but they are typically used for different purposes. Qiskit's Aer module is particularly suited to noisy circuit simulation, including channel loss and measurement error, whereas PennyLane is often used in hybrid quantum-classical workflows and optimisation focused studies. For this project, Qiskit combined with the Aer module is the more appropriate choice because the focus is on protocol-level benchmarking under controlled noise assumptions rather than variational or hybrid optimisation.

\section{Research Gaps}
The existing \gls{qkd} simulation literature reveals several gaps that this project seeks to address. A common limitation is that many studies evaluate \gls{bb84}, \gls{b92}, or \gls{e91} in isolation, making it difficult to compare protocol behaviour under common assumptions and metrics \cite{Ma2005, Tamaki2014}. Furthermore, published simulations often differ in their treatment of noise, eavesdropping, and implementation constraints, meaning that results are not always directly comparable across protocol families.

This project therefore aims to develop a unified simulation framework to benchmark \gls{bb84}, \gls{b92}, and \gls{e91} within the same environment and under consistent assumptions. By comparing \gls{pandm} and \gls{eb} schemes using the same metrics, the study aims to produce a more controlled basis for evaluating efficiency, and practical suitability.

A further limitation in the literature is that many simulation studies prioritise theoretical metrics without sufficiently connecting them to implementation realism. Using the Qiskit Aer module enables controlled noise modelling and protocol-level analysis, which makes the study more relevant to both research comparison and practical prototyping \cite{QiskitAer}.



%================================================================================
%                  Theoretical BG
%================================================================================
\chapter{Theoretical Background}
%================================================================================

\section{Quantum Mechanical Foundations}

To achieve \gls{itsec}, every \gls{qkd} protocol rests on a small set of quantum mechanical principles. This section introduces the core principles exploited by the \gls{bb84}, \gls{b92}, and \gls{e91} protocols, following the treatment by Nielsen and Chuang \cite{QKDbible}.

\subsection{Qubits \& Superposition}
While a classical bit takes a definite value of 0 or 1, a \gls{qubit} can exist in a linear combination of both. Using Dirac notation, the general single-\gls{qubit} state is written as
%
\begin{equation}
\label{eq:qubit}
\ket{\psi} = \alpha\ket{0} + \beta\ket{1},
\end{equation}
%
where $\alpha$ and $\beta$ are complex probability amplitudes satisfying $|\alpha|^{2}+|\beta|^{2}=1$ \cite{QKDbible}. The states $\ket{0}$ and $\ket{1}$ form the rectilinear ($Z$) basis. A \gls{qubit} can also be expressed in the diagonal ($X$) basis, defined as
%
\begin{equation}
\label{eq:diag-basis}
\ket{+} = \frac{\ket{0}+\ket{1}}{\sqrt{2}}, \qquad
\ket{-} = \frac{\ket{0}-\ket{1}}{\sqrt{2}}.
\end{equation}
%
In physical \gls{qkd} systems, \glspl{qubit} are typically encoded in photon polarisation, where the $Z$ basis corresponds to rectilinear states and the $X$ basis to $\pm 45^{\circ}$ diagonal states \cite{BB84}. 

\subsection{Measurement Disturbance}
Superposition becomes a security asset because measuring a \gls{qubit} irreversibly forces a basis choice. Measuring in a compatible basis yields the correct outcome, but measuring in an incompatible basis destroys the original state, collapsing it into a random eigenstate with equal probability \cite{QKDbible}. This basis dependence underpins \gls{sifting} in \gls{bb84} and \gls{b92}, where only compatible-basis measurements contribute to the final key \cite{BB84, B92}.

\subsection{No-Cloning Theorem}
The no-cloning theorem \cite{WZ1982, Dieks1982} proves that no physical process can produce a perfect copy of an arbitrary unknown quantum state. This theorem provides the fundamental guarantee that an eavesdropper cannot silently duplicate quantum transmissions, granting \gls{qkd} its information-theoretic security \cite{ekertQI}.

\subsection{Heisenberg's Uncertainty}
Closely related to measurement disturbance, Heisenberg's uncertainty principle \cite{Heisenberg1927} states that non-commuting observables (such as $Z$ and $X$ basis polarisations) cannot be simultaneously known with precision \cite{QKDbible}. In \gls{qkd}, this forces a physical trade-off: any information gained by an eavesdropper in one basis inevitably introduces detectable errors in the conjugate basis, leaving a measurable signature of interception.

\subsection{Quantum Entanglement \& Bell's Theorem}
Two or more \glspl{qubit} can exist in an entangled state that cannot be determined independently. The Bell state (\gls{epr} pair) \cite{EPR1935} is defined as:
%
\begin{equation}
\label{eq:bell-state}
\ket{\Phi^{+}} = \frac{\ket{00}+\ket{11}}{\sqrt{2}}.
\end{equation}
%
Measuring one \gls{qubit} instantaneously projects the other into a correlated state, regardless of spatial separation \cite{QKDbible}. Bell \cite{Bell1964} formalised this non-local correlation, and the \gls{chsh} inequality \cite{CHSH1969} provides an experimentally testable bound:
%
\begin{equation}
\label{eq:chsh}
|S| = |E(a,b) - E(a,b') + E(a',b) + E(a',b')| \leq 2.
\end{equation}
%
Quantum mechanics predicts a maximum violation of $|S| = 2\sqrt{2} \approx 2.828$ \cite{Aspect1982}. Ekert's \gls{e91} protocol \cite{E91} leverages this violation directly as a built-in eavesdropping test to verify channel security.


\section{Classical Post-Processing Foundations}
\label{sec:postfoundations}

Evaluating \gls{qkd} performance requires translating quantum measurement outcomes into three classical metrics: the \gls{qber}, Mutual Information, and the \gls{skr}. These form an analytical chain through physical degradation, channel capacity, and cryptographic yield, applied uniformly across all three protocols in \Cref{sec:postprocessing}.


\subsection{Mutual Information}
\label{subsec:mutualinfobg}

Because depolarising noise applies errors symmetrically, the sifted key relationship can be modelled as a Binary Symmetric Channel (\gls{bsc}) with crossover probability equal to the \gls{qber} ($E$). The Mutual Information $I_{AB}$, as described by Johnston \cite{mutualInfo}, represents the capacity of this channel:

\begin{equation}
    I_{AB} = 1 - H(E)
    \label{eq:mi_theory}
\end{equation}

where $H(x) = -x\log_2(x) - (1-x)\log_2(1-x)$ is the binary Shannon entropy. A secure key exists only if Alice and Bob share more information than Eve: $I_{AB} > I_{AE}$.

\subsection{Devetak-Winter Bound}
\label{subsec:dwboundbg}

The Devetak-Winter bound \cite{DWbound2005} formalises the maximum achievable \gls{skr} fraction $K$:

\begin{equation}
    K = I_{AB} - \max I_{AE}
    \label{eq:dw_theory}
\end{equation}

Because Alice and Bob cannot distinguish environmental noise from eavesdropping, information-theoretic security requires that all channel errors are attributed to an adversary deploying an optimal attack. The form of $\max I_{AE}$ therefore depends on the optimal attack model for each protocol, as applied in \Cref{subsec:DWsec}.



%================================================================================
%                  Methodology
%================================================================================
\chapter{Methodology}
%================================================================================
\section{Simulation Architecture}
\label{sec:simarch}

\subsection{Software Development Methodology}
\label{subsec:dev}

The simulation was developed using an iterative and incremental software engineering methodology in Python, utilising the Qiskit \gls{sdk}, Matplotlib and NumPy libraries, Github, and MATLAB to perform 3D surface plotting. This approach was chosen to mitigate development risk and ensure the mathematical validity of each \gls{qkd} protocol before introducing broader system features.

\subsection{Modular Package Design}
\label{subsec:moddesign}

Implemented as a modular Python package (\texttt{QKDSimulator}), the system is functionally segmented, as depicted in \Cref{fig:systemfiles}, enabling independent protocol operation and reproducible experiment configuration via decoupled parameters.

\begin{figure}[htbp]
    \centering
    \includegraphics[width=1\linewidth]{images.methodology/system overview.png}
    \caption{Modularised System Architecture.}
    \label{fig:systemfiles}
\end{figure}

The architecture employs an \gls{oop} paradigm organised around a parent class, \texttt{QKDProtocol}, which defines the common interface shared by all three protocols. Each child class (\gls{bb84}, \gls{b92}, \gls{e91}) overrides the parent with independent operation logic while inheriting common post-processing functions. A dataclass (\texttt{QKDResult}) encapsulates the output runs (sifted keys, \gls{qber}, and derived metrics), ensuring protocol data is processed interchangeably by the benchmark module.

\newpage
\subsection{Experiment Configuration System}
\label{subsec:config}

During runtime, the selected configuration file is parsed into a \texttt{SimConfig} dataclass located in \texttt{config.py}, which validates the requested execution mode against the four supported modes: \texttt{single} (a single protocol test trial with fixed parameters), \texttt{sweep} (a parametric evaluation of channel Noise Strength), \texttt{eve\_sweep} (a parametric evaluation of eavesdropper interception rates), and \texttt{protocol\_comparison} (execution of multiple protocols under identical noise conditions). An example configuration snippet is provided in \Cref{lst:yaml-config}.

\begin{lstlisting}[language=Python, caption={Example YAML Configuration for a BB84 Noise Sweep.}, label={lst:yaml-config}]
mode: sweep
protocol: bb84
n_qubits: 1000
n_trials: 20
noise_type: depolarizing
noise_min: 0.0
noise_max: 0.30
noise_steps: 30
output_dir: ./results/bb84_depolarizing
save_plots: true
\end{lstlisting}

\subsection{System Information Flow}
\label{subsec:infoflow}

Building upon the configuration system, the execution of a simulation follows a deterministic, sequential information flow across four phases: Preparation, Measurement, Post-processing, and Plotting, as illustrated in \Cref{fig:pipeline}.

To improve scalability across large-scale parameter sweeps, the measurement phase employs batch processing. Rather than dispatching the \texttt{AerSimulator} iteratively per \gls{qubit}, all prepared Qiskit circuits are aggregated and submitted as a single batch job, as shown in \Cref{lst:individual} and \Cref{lst:batch}.

\begin{lstlisting}[language=Python,
  caption={Individual Circuit Processing.},
  label={lst:individual}]
def _bob_measure(self, circuits: List[QuantumCircuit]):
    """Bob picks random bases and measures each qubit sequentially."""
    self._bob_bases   = np.random.randint(0, 2, size=self.n_qubits)
    self._bob_results = np.zeros(self.n_qubits, dtype=int)

    for i, qc in enumerate(circuits):
        if self._bob_bases[i] == 1:
            qc.h(0)
        qc.measure(0, 0)

        job    = self.backend.run([qc], shots=1, memory=True)
        result = job.result()

        self._bob_results[i] = int(result.get_memory(0)[0])
\end{lstlisting}

\begin{lstlisting}[language=Python,
  caption={Improved Batch Circuit Processing.},
  label={lst:batch}]
def _bob_measure(self, circuits: List[QuantumCircuit]):
    """Bob picks random bases and measures each qubit as a batch."""
    self._bob_bases   = np.random.randint(0, 2, size=self.n_qubits)
    self._bob_results = np.zeros(self.n_qubits, dtype=int)

    for i, qc in enumerate(circuits):
        if self._bob_bases[i] == 1:
            qc.h(0)
        qc.measure(0, 0)

    job    = self.backend.run(circuits, shots=1, memory=True)
    result = job.result()
    for i in range(len(circuits)):
        self._bob_results[i] = int(result.get_memory(i)[0])
\end{lstlisting}

If specified in the configuration file, an \texttt{EveInterceptor} object is initialised during the Preparation phase. Eve's interception acts as a modular interrupt within the standard \gls{pandm} pathway, enabling the threat model to be integrated without alteration to Bob's fundamental measurement logic.

For Experiments 5 \& 6, which evaluate the multivariable dynamics of
\gls{bb84} and \gls{b92}, trial data is collected at runtime via the
\texttt{SurfaceBenchMarkData} object and exported as a structured
\texttt{.csv} file for subsequent MATLAB surface rendering.

\begin{figure}[htbp]
    \centering
    \includegraphics[width=1\linewidth]{images.methodology/information flow.png}
    \caption{System Information Flow}
    \label{fig:pipeline}
\end{figure}



\newpage
\section{Protocol Implementation}
\label{sec:protocolimplementation}

\Cref{tab:protocol-comparison} defines the qualitative characteristics of each protocol that served as the implementation reference.

% qualitative table
\begin{table}[htbp]
\centering
\caption{Qualitative Comparison of BB84, B92, and E91.}
\label{tab:protocol-comparison}
\footnotesize
\begin{tabularx}{\linewidth}{@{}l X X X@{}}
\toprule
\textbf{Feature} & \textbf{BB84} & \textbf{B92} & \textbf{E91} \\
\midrule
Paradigm & \gls{pandm} & \gls{pandm} & \gls{eb} \\
Encoding States & 4 (Two Bases) & 2 Non-Orthogonal & Bell Pairs, 3 Angles \\
Alice's Bases / Angles & $\{Z, X\}$ & $\{\ket{0}, \ket{+}\}$ & $\{0, \pi/4, \pi/2\}$ \\
Bob's Bases / Angles & $\{Z, X\}$ & $\{Z, X\}$ & $\{\pi/4, \pi/2, 3\pi/4\}$ \\
\gls{sifting} Efficiency & $\sim$50\% & $\sim$25\% & $\sim$22.2\% \\
Security Threshold & $\sim$11\% & $\sim$6.5\% & $\sim$11\% \\
Eve Detection & \gls{qber} & \gls{qber} & \gls{chsh} $+$ \gls{qber} \\
Full Intercept-Resend \gls{qber} & $\sim$25\% & $\sim$38\% & Degraded $|S|$ \\
Known Vulnerabilities & PNS & USD & Entanglement Loss \\
Hardware Complexity & Moderate & Lowest & Highest \\
\bottomrule
\end{tabularx}
\end{table}

\subsection{BB84}
\label{subsec:bb84}

\gls{bb84} distributes key material by encoding bits across two mutually unbiased bases (rectilinear and diagonal), with Alice selecting a random bit value and basis per \gls{qubit} to produce one of the four states in \Cref{eq:bb84-states} \cite{BB84}, illustrated in \Cref{fig:BB84operation}.

%
\begin{equation}
  \ket{0},\quad \ket{1},\quad
  \ket{+} = \frac{\ket{0}+\ket{1}}{\sqrt{2}},\quad
  \ket{-} = \frac{\ket{0}-\ket{1}}{\sqrt{2}}.
  \label{eq:bb84-states}
\end{equation}
%

\begin{figure}[htbp]
    \centering
    \includegraphics[width=1\linewidth]{images.methodology/bb84processTNR.png}
    \caption{BB84 Protocol Operation Model.}
    \label{fig:BB84operation}
\end{figure}

\newpage
Following transmission, Alice and Bob compare basis choices over the authenticated \gls{cc}, discarding mismatched rounds; under ideal conditions, $\sim\!50\%$ of \glspl{qubit} survive \gls{sifting}.

As illustrated in \Cref{fig:bb84QC}, state encoding is conditional: an X gate and optional Hadamard rotate each \gls{qubit} into the required basis. Critically, an Identity gate (\texttt{qc.id(0)}) is appended immediately after preparation. This acts as the \gls{qchan}, providing a point to inject noise and eavesdrop without altering the logical quantum state.

Bob measures each incoming \gls{qubit} in a randomly selected basis; following circuit execution, basis arrays are compared and the $\sim\!50\%$ sifted key is passed to the post-processing pipeline.

\begin{figure}[htbp]
    \centering
    \includegraphics[width=1\linewidth]{images.methodology/bb84QC.png}
    \caption{Qiskit Circuit Model for BB84.}
    \label{fig:bb84QC}
\end{figure}

\newpage
\subsection{B92}
\label{subsec:b92}

\gls{b92} reduces the \gls{bb84} encoding to two non-orthogonal states, relying on the principle that no measurement can perfectly distinguish them to guarantee security \cite{ekertQI}. Alice encodes each bit using the scheme in \Cref{eq:b92-states}, with the $45^\circ$ angular separation between states confirmed in \Cref{fig:b92basis}.

%
\begin{equation}
  \ket{\varphi} =
  \begin{cases}
    \ket{0}  & \text{if } a = 0, \\[4pt]
    \ket{+}  & \text{if } a = 1.
  \end{cases}
  \label{eq:b92-states}
\end{equation}
%

\begin{figure}[htbp]
    \centering
    \includegraphics[width=0.7\linewidth]{images.methodology/b92basis.png}
    \caption{B92 Preparation \& Measurement Basis Angles.}
    \label{fig:b92basis}
\end{figure}

Bob's measurement logic, detailed in \Cref{fig:b92flowchart}, is the most distinctive feature of \gls{b92}. A measurement outcome is only conclusive when it logically contradicts one of Alice's two possible preparations: an outcome of ``\(1\)'' in the $Z$-basis rules out $\ket{0}$, whilst an outcome of ``\(1\)" in the $X$-basis rules out $\ket{+}$.

Unlike \gls{bb84}, Alice and Bob do not exchange basis information during \gls{sifting}. Bob publicly announces only the indices of conclusive events (outcome ``\(1\)''); all inconclusive rounds are discarded. For each conclusive round, Bob's sifted key bit is the logical complement of his measurement basis, $1 - b_i$.

Because \gls{b92} relies on only two states and discards ambiguous outcomes, its \gls{sifting} efficiency is fundamentally lower than that of \gls{bb84}, with a theoretical efficiency of $\sim\!25\%$ under ideal conditions. In practical deployment, this leaves the protocol more susceptible to channel noise, requiring longer transmission sequences to generate an equivalent key length.

\begin{figure}[htbp]
    \centering
    \includegraphics[width=1\linewidth]{images.methodology/b92 flowchart.png}
    \caption{B92 Algorithm Flowchart.}
    \label{fig:b92flowchart}
\end{figure}



\newpage
\subsection{E91}
\label{subsec:e91}

Unlike \gls{bb84} and \gls{b92}, \gls{e91} establishes a shared key by exploiting quantum \gls{entanglement}, with security verified through violations of the \gls{chsh} inequality rather than the no-cloning theorem alone \cite{E91}. The protocol operates on maximally entangled \gls{qubit} pairs in the singlet Bell state:

\begin{equation}
\ket{\Psi^-} = \frac{\ket{01} - \ket{10}}{\sqrt{2}}
\label{eq:e91_singlet}
\end{equation}

Security is verified via the \gls{chsh} parameter $S$, computed from correlation coefficients across non-aligned angle pairs:

\begin{equation}
S = E(a_1, b_1) - E(a_1, b_3) + E(a_3, b_1) + E(a_3, b_3)
\label{eq:chsh_parameter}
\end{equation}

As illustrated in \Cref{fig:e91angles}, the measurement angles are chosen to satisfy the angular separations required by the \gls{chsh} test and to yield the two aligned pairs used for key generation, with the corresponding circuit shown in \Cref{fig:e91circuit}.

\begin{figure}[htbp]
    \centering
    \includegraphics[width=0.65\linewidth]{images.methodology/e91 XZ plane angles.png}
    \caption{E91 XZ-Plane Cross-section of Measurement Angles.}
    \label{fig:e91angles}
\end{figure}

\begin{figure}[htbp]
    \centering
    \includegraphics[width=0.8\linewidth]{images.methodology/e91_circuit.png}
    \caption{E91 Quantum Circuit.}
    \label{fig:e91circuit}
\end{figure}

\newpage
%\subsubsection*{State Preparation}

Since the standard Hadamard $\rightarrow$ CNOT sequence in Qiskit produces $\ket{\Phi^+} = \frac{1}{\sqrt{2}}(\ket{00} + \ket{11})$ rather than the required $\ket{\Psi^-}$ singlet, the system logic appends an X gate followed by a Z gate to the first \gls{qubit} after the entangling operation, as described in \Cref{alg:e91_prep} and illustrated in \Cref{fig:e91circuit}.

\begin{algorithm}[htbp]
    \caption{E91 Singlet State Preparation.}\label{alg:e91_prep}
    \KwIn{Two-\gls{qubit} quantum circuit \texttt{qc}, \gls{qubit} indices \texttt{q[0]} (Alice), \texttt{q[1]} (Bob)}
    \KwOut{Circuit \texttt{qc} encoding $\ket{\Psi^-}$ across \texttt{q[0]} and \texttt{q[1]}}
    Apply Hadamard gate $H$ to \texttt{q[0]}\; \Comment{Initialise \texttt{q[0]} into an equal superposition}
    Apply CNOT gate with \texttt{q[0]} as control, \texttt{q[1]} as target\; \Comment{Entangle \texttt{q[0]} and \texttt{q[1]} to produce $\ket{\Phi^+}$}
    Apply X gate to \texttt{q[0]}\; \Comment{Flip computational basis of \texttt{q[0]} -- transforms $\ket{\Phi^+} \rightarrow \ket{\Psi^+}$}
    Apply Z gate to \texttt{q[0]}\; \Comment{Apply relative phase flip to \texttt{q[0]} -- transforms $\ket{\Psi^+} \rightarrow \ket{\Psi^-}$}
    \Return{\texttt{qc}}
\end{algorithm}

This four-step sequence is vital because $\ket{\Psi^-}$ possesses a relative phase of $-1$ between its $\ket{01}$ and $\ket{10}$ components. The Z gate introduces this phase relationship, whilst the X gate performs the computational basis flip required to produce anti-correlated rather than correlated outcomes \cite{QiskitAer}. Selecting $\ket{\Psi^-}$ over $\ket{\Phi^+}$ is integral to the angular measurement decoding logic, as it ensures anti-correlation.

Measurement along an arbitrary angle $\theta$ in the XZ-plane is simulated by applying $R_y(-\theta)$ to the relevant \gls{qubit} prior to $Z$-basis projection, making the computational-basis measurement equivalent to a physical polarisation measurement at angle $\theta$ \cite{e91Qiskit}. \Cref{alg:e91_meas} presents this procedure.

Following measurement, outcomes are categorised against the $3 \times 3$ \gls{sifting} table indexed by the nine possible Alice-Bob angle-pair combinations, as shown in \Cref{tab:e91_sift}.

Of the nine combinations, two pairs are designated for key generation, four for \gls{chsh} estimation, and three are discarded. For key pairs, Bob logically inverts his outcomes ($1 - b_i$) to reconcile the singlet anti-correlation and align the shared key. Since each combination is equally probable, the theoretical \gls{sifting} efficiency is $2/9 \approx 22.2\%$, a trade-off inherent to the protocol's security verification overhead.

\newpage
\begin{table}[htbp]
\centering
\caption{E91 Sifting Table. Key Pairs Contribute to the Raw Sifted Key; CHSH Pairs Estimate the Bell Parameter $S$; Discard Pairs are Rejected.}
\label{tab:e91_sift}
\footnotesize
\begin{tabular}{lccc}
\toprule
 & \textbf{Bob: $45^\circ$} & \textbf{Bob: $90^\circ$} & \textbf{Bob: $135^\circ$} \\
\midrule
\textbf{Alice: $0^\circ$}  & \gls{chsh} & \gls{chsh} & Discard \\
\textbf{Alice: $45^\circ$} & Key        & \gls{chsh} & \gls{chsh} \\
\textbf{Alice: $90^\circ$} & Discard    & Key        & Discard \\
\bottomrule
\end{tabular}
\end{table}

\begin{algorithm}[htbp]
    \caption{E91 Angular Measurement.}\label{alg:e91_meas}
    \KwIn{Circuit \texttt{qc} with $\ket{\Psi^-}$ prepared, angle sets $A_{\text{alice}} = \{0^\circ, 45^\circ, 90^\circ\}$, $A_{\text{bob}} = \{45^\circ, 90^\circ, 135^\circ\}$, shot count \texttt{n\_shots}}
    \KwOut{Measurement result list \texttt{outcomes[]}}
    \For{$i \leftarrow 1$ \KwTo \texttt{n\_shots}}{
        Sample angle $a$ randomly from $A_{\text{alice}}$\; \Comment{Alice randomly selects her measurement angle}
        Sample angle $b$ randomly from $A_{\text{bob}}$\; \Comment{Bob randomly selects his measurement angle}
        Apply $R_y(-a)$ to \texttt{q[0]}\; \Comment{Rotate Alice's state by $-a$ -- aligns measurement with chosen basis}
        Apply $R_y(-b)$ to \texttt{q[1]}\; \Comment{Rotate Bob's state by $-b$ -- aligns measurement with chosen basis}
        Measure \texttt{q[0]} $\rightarrow$ \texttt{c[0]}, \texttt{q[1]} $\rightarrow$ \texttt{c[1]}\; \Comment{Execute projective measurement in the $Z$-basis}
        Append $(a, b, \texttt{c[0]}, \texttt{c[1]})$ to \texttt{outcomes[]}\;
    }
    \Return{\texttt{outcomes[]}}
\end{algorithm}

\newpage
\section{QKD Network Topology}
\label{sec:networktopology}

The \gls{pandm} and \gls{eb} topologies impose different simulation requirements for state distribution, channel noise, and eavesdropping. Across both, the \gls{cc} is assumed perfectly authenticated and impervious to attack, as \gls{qkd} depends on this pre-established classical authenticity during post-processing.

\subsection{Prepare \& Measure Topology}
\label{subsec:pmtopology}

As depicted in \Cref{fig:pmtopology}, only one travelling \gls{qubit} is modelled per round; channel interference is therefore inserted at a single point between preparation and measurement.

\begin{figure}[htbp]
    \centering
    \includegraphics[width=1\linewidth]{images.methodology/QuantumAbstractTNR.png}
    \caption{Abstract QKD Network Topology of Prepare \& Measure Protocols.}
    \label{fig:pmtopology}
\end{figure}


\newpage
\subsection{Entanglement-Based Topology}
\label{subsec:ebtopology}

As shown in \Cref{fig:ebtopology}, the \gls{eb} topology replaces the unidirectional channel with a shared \gls{epr} source; security is inferred from measurement correlations rather than state disturbance, and channel interference acts on both \glspl{qubit} simultaneously \cite{ekertQI}.

\begin{figure}[htbp]
    \centering
    \includegraphics[width=0.8\linewidth]{images.methodology/EBtopology.png}
    \caption{Abstract QKD Network Topology of Entanglement-Based Protocols.}
    \label{fig:ebtopology}
\end{figure}

\newpage
\section{Quantum Channel Interference}
\label{sec:channel_interference}

Two interference sources are modelled: depolarising noise (environmental decoherence) and an intercept-resend attack (adversarial interception).

\subsection{Depolarising Noise}
\label{subsec:noise}

The depolarising channel model is adopted as a standard worst-case representation of environmental decoherence \cite{QKDbible}, replacing the quantum state with a mixed state with probability $p$ and leaving it error-free with probability $1-p$. It maps an arbitrary single-\gls{qubit} state $\rho$ by applying a random Pauli error with total probability $p$:

\begin{equation}
\label{eq:depolarising}
\mathcal{E}_{dep}(\rho) = (1-p)\rho + \frac{p}{3}(X\rho X + Y\rho Y + Z\rho Z)
\end{equation}

Consequently, as the Noise Strength parameter $p$ increases, the probability of Bob measuring an erroneous outcome grows, directly increasing the \gls{qber}.

To implement this without hardcoding errors into the protocol logic, the platform exploits Qiskit's \texttt{AerSimulator}. However, a challenge arises: by default Qiskit noise models apply errors to specific gate types across the entire circuit. Assigning a depolarising error to a standard operation such as the X or Hadamard gate would corrupt the \gls{qubit} during Alice's preparation and Bob's measurement phases, simulating an imperfect source or receiver rather than a noisy \gls{qchan}.

Resolving this required binding noise exclusively to the Identity gate (\texttt{id}), which is appended to each circuit at the transmission point as described in \Cref{sec:protocolimplementation}. As detailed in \Cref{alg:noise_model}, the depolarising error targets only \texttt{['id']} gates, ensuring preparation and measurement remain mathematically perfect and that any resultant \gls{qber} is an isolated product of the simulated \gls{qchan}.

\begin{algorithm}[htbp]
    \caption{AerSimulator Backend Generation for Channel Noise.}\label{alg:noise_model}
    \KwIn{\texttt{noise\_type} (e.g. 'depolarizing'), \texttt{strength} $p$}
    \KwOut{Configured \texttt{AerSimulator} backend}
    \If{\texttt{noise\_type} $=$ 'none'}{
        \Return{default \texttt{AerSimulator()}} \Comment{Ideal channel -- no noise applied}
    }
    \texttt{noise\_model} $\leftarrow$ empty \texttt{NoiseModel()}\;
    \texttt{error} $\leftarrow$ generate \texttt{depolarizing\_error} with probability \texttt{strength}\; \Comment{Define depolarising error parameterized by $p$}
    Add \texttt{error} to \texttt{noise\_model}, targeting only \texttt{['id']} gates\; \Comment{Apply noise model exclusively to identity operations}
    \Return{\texttt{AerSimulator(noise\_model=noise\_model)}}
\end{algorithm}


\newpage
\subsection{Eavesdropping}
\label{subsec:eavesdropping}

In the intercept-resend attack, Eve intercepts each \gls{qubit}, measures in a randomly guessed basis, and resends a replacement state encoded with her result. Incorrect basis guesses introduce detectable errors, elevating the \gls{qber} above the channel noise floor.

Within the simulator, this is encapsulated in the \texttt{EveInterceptor} class, with attack intensity governed by a configurable \texttt{interception\_rate} parameter ($0.0$ \text{to} $1.0$) that sets the per-\gls{qubit} interception probability. \Cref{alg:eve_intercept} details the logical execution of the attack.

\begin{algorithm}[htbp]
    \caption{Intercept-Resend Eavesdropper Logic.}\label{alg:eve_intercept}
    \KwIn{List of Alice's prepared \texttt{circuits}, \texttt{interception\_rate} parameter}
    \KwOut{Modified list of circuits to be received by Bob}
    \For{each \texttt{qc} in \texttt{circuits}}{
        \If{\text{random\_value()} $< \texttt{interception\_rate}$}{
            \texttt{eve\_basis} $\leftarrow$ randomly select measurement basis ($0$ or $1$)\;
            Measure \texttt{qc} in \texttt{eve\_basis}, yielding classical bit \texttt{eve\_result}\; \Comment{Collapse the intercepted state to yield a classical bit}
            \texttt{new\_qc} $\leftarrow$ initialise empty quantum circuit\;
            Encode \texttt{eve\_result} into \texttt{new\_qc} using \texttt{eve\_basis}\; \Comment{Prepare new quantum state for resending}
            Append \texttt{id} channel marker to \texttt{new\_qc}\; \Comment{Insert identity gate to facilitate subsequent channel noise modelling}
            Replace original \texttt{qc} with \texttt{new\_qc}\;
        }
    }
    \Return{modified \texttt{circuits}}
\end{algorithm}

Critically, the replacement circuit retains the \texttt{id} gate, preserving the ability to model channel noise on the resend leg and accurately reflecting a noisy \gls{qchan} under active attack.

\subsubsection{E91 Limitations}

The \texttt{EveInterceptor} module is constrained to the \gls{pandm} schemes. Because \gls{e91} relies on distributed \gls{entanglement}, an intercept-resend attack on one arm of the \gls{epr} pair breaks the singlet state entirely, requiring Eve to instead execute \gls{entanglement}-swapping operations or a Bell-state disruption strategy \cite{E91}, both are threats beyond the scope of this project.



\newpage
\section{Classical Post-Processing Pipeline}
\label{sec:postprocessing}

Reconciliation and privacy amplification are excluded due to computational overhead, with the pipeline instead computing QBER, Mutual Information, and SKR under the Shannon limit \cite{Shannon1948}.


\subsection{QBER Estimation}
\label{subsec:QBER}

The \gls{qber} serves as the primary security parameter for \gls{pandm} protocols, defined as the ratio of errors ($N_{errors}$) to the total length of the sifted key ($N_{sifted}$). Within the simulation, Alice and Bob's bit arrays are directly accessible in memory, hence, the post-processing pipeline bypasses sacrificial sampling and computes the \gls{qber} across the entire sifted array to maximise precision. Once the valid \texttt{sifted\_key\_alice} and \texttt{sifted\_key\_bob} arrays are isolated, the calculation is executed using vectorised NumPy operations to ensure high computational throughput over large-scale samples, as demonstrated in \Cref{lst:qber_calc}.

\begin{lstlisting}[language=Python, caption={Vectorised QBER Calculation Logic.}, label={lst:qber_calc}]
sifted_length = len(sifted_key_alice)
if sifted_length > 0:
    errors = np.sum(sifted_key_alice != sifted_key_bob)
    qber = errors / sifted_length
else:
    qber = 0.0
\end{lstlisting}

\subsection{Mutual Information}
\label{subsec:mutualinfo}

As established in \Cref{sec:postfoundations}, the sifted key relationship is modelled as a \gls{bsc}, giving $I_{AB} = 1 - H(E)$, as described by Johnston \cite{mutualInfo}. A secure key requires $I_{AB} > I_{AE}$; Eve's presence under the intercept-resend model inherently elevates $E$, shifting this balance until the intersection $I_{AB} = I_{AE}$ defines the security threshold. For \gls{bb84}, this occurs at $E \approx 11\%$, above which \gls{amplification} cannot distil a secure key and transmission is aborted
\cite{ShorPreskill2000}.

\subsection{Security Verification: CHSH Parameter}
\label{subsec:CHSH}

As established in \Cref{subsec:e91}, the \gls{chsh} Bell parameter $S$ is calculated from four specific correlation coefficients $E(a_i, b_j)$, corresponding to the angle-pairs designated for \gls{chsh} estimation in \Cref{tab:e91_sift}. To calculate the $S$-parameter, the pipeline takes measurement outcomes corresponding to the \gls{chsh} subset. For each of the four angle pairings, raw binary outcomes are mapped to parity values ($+1$ or $-1$), and the correlation coefficient is computed as the mean of the product of Alice and Bob's parity arrays. \Cref{lst:chsh_calc} demonstrates the vectorised loop used to execute this extraction.

\begin{lstlisting}[language=Python, caption={Vectorised Extraction of CHSH Correlation Coefficients.}, label={lst:chsh_calc}]
# _CHSH_S_TERMS = [((0, 0), +1), ((0, 2), -1), ((2, 0), +1), ((2, 2), +1)]
s_value = 0.0
for (ai, bi), sign in self._CHSH_S_TERMS:
    # Isolate the subset of measurements for the angle pair
    sub_mask = (self._alice_angles == ai) & (self._bob_angles == bi)
    
    if np.sum(sub_mask) > 0:
        # Map binary (0,1) outcomes to parity (+1,-1)
        a_pm = 1 - 2 * self._alice_outcomes[sub_mask].astype(int)
        b_pm = 1 - 2 * self._bob_outcomes[sub_mask].astype(int)
        
        # Calculate expected correlation and add to total S parameter
        corr = float(np.mean(a_pm * b_pm))
        s_value += sign * corr
\end{lstlisting}

Aggregating these terms produces a final \texttt{s\_value} for each \gls{e91} run, which serves as the security flag for key distillation.

\newpage
\subsection{Devetak-Winter Secure Key Rate Formula}
\label{subsec:DWsec}

Applying the Devetak-Winter bound established in \Cref{sec:postfoundations}, the system assumes operation at the Shannon capacity limit throughout. As modelling \gls{reconciliation} is outside the scope of this simulator, error correction is treated as perfectly efficient.

Under worst-case attribution (\Cref{sec:postfoundations}), the pipeline applies a protocol-specific $\max I_{AE}$ based on the optimal attack model for each protocol, determined by its underlying physics:

\begin{itemize}
    \item \textbf{BB84:} Eve must guess between two mutually unbiased bases. Thus, under the Shannon limit, the Devetak-Winter bound reduces to $K = 1 - 2H(E)$, mirroring the GLLP limit \cite{GLLP2004}.
    \item \textbf{B92:} Eve must perform a USD attack. For states separated by $\theta = 45^\circ$, the penalty is: $I_{AE} = H\left(\frac{E}{1 - \cos 45^\circ}\right)$.
    \item \textbf{E91:} Eve's information is bounded by the degradation of the \gls{chsh} parameter $S$: $I_{AE} = H\left(\frac{1 + \sqrt{(S/2)^2 - 1}}{2}\right)$, adapted from Acin \textit{et al.} \cite{Acin2007}.
\end{itemize}

\Cref{lst:dw_bound} details this implementation.

\begin{lstlisting}[language=Python, caption={Dynamic Application of the Devetak-Winter Bound.}, label={lst:dw_bound}]
@property
def secure_key_rate(self) -> float:
    """Calculates the Devetak-Winter SKR assuming the Shannon limit"""
    
    if self.sifted_length == 0:
        return 0.0

    # Information securely shared between Alice and Bob
    i_ab = self.mutual_information
    
    # Calculate Eve's penalty based on the specific protocol physics
    if self.protocol_name == "BB84":
        if self.qber >= 0.110:
            return 0.0
        i_ae = self.key_rate * self._binary_entropy(self.qber)
        
    elif self.protocol_name == "B92":
        if self.qber >= 0.065:
            return 0.0
        # Penalty assuming a USD attack
        denominator = 1.0 - (1.0 / np.sqrt(2.0))
        i_ae = self.key_rate * self._binary_entropy(self.qber / denominator)
        
    elif self.protocol_name == "E91":
        s_val = getattr(self, 'abs_s', 0.0)
        if s_val <= 2.0:
            return 0.0
        # Penalty derived from the CHSH Bell parameter degradation
        penalty = (1.0 + np.sqrt(max(0.0, (s_val / 2.0)**2 - 1.0))) / 2.0
        i_ae = self.key_rate * self._binary_entropy(penalty)
        
    else:
        return 0.0

    return max(0.0, i_ab - i_ae)
\end{lstlisting}

If $K \leq 0$, the channel is either too noisy or compromised and the transmission is aborted.

\newpage
\section{Experimental \& Statistical Methods}
\label{sec:experimental_methods}

To evaluate the protocols under varying channel conditions, the platform executes a series of benchmarking sweeps using a Monte Carlo approach to extract reliable performance metrics.

\subsection{Central Limit Theorem and Monte Carlo Simulation}
\label{subsec:clt_monte_carlo}

The experiments execute $N_{trials}$ independent runs at every defined parameter point, such as a specific noise strength ($p$) or interception rate. By following the Central Limit Theorem \cite{montgomery2018}, the platform aggregates these independent trials to calculate the sample mean ($\mu$) and the standard error of the mean (SEM) for each metric:

\begin{equation}
\text{SEM} = \frac{\sigma}{\sqrt{N_{trials}}}
\label{eq:sem}
\end{equation}

where $\sigma$ is the sample standard deviation across the trials. This enables the results to represent statistical protocol behaviour, rather than anomalous measurements.

\subsection{Protocol Performance Metrics}
\label{subsec:performance_metrics}

Three independent variables exist: a baseline control channel, dynamic depolarising strength ($0.0 \leq p \leq 0.5$), and a dynamic interception rate ($0.0 \leq \text{Eve Rate} \leq 1.0$). As for dependent variables; \gls{qber}, $I(A;B)$, $I(A;E)$, Devetak-Winter \gls{skr} ($K$), and $|S|$ are computed as per \Cref{sec:postprocessing}. By correlating these metrics across identical channel noise profiles, the simulation framework provides a uniform basis for comparing protocol performance.


\section{Data Visualisation}
\label{sec:dataviz}

While simulation and statistical reduction are performed within Python, presenting multi-dimensional results requires dedicated rendering. Matplotlib's 3D backend is constrained by its rasterised rendering, preventing the precise control over surfaces required for high-quality figures. The visualisation workflow is therefore partitioned, with Python executing the simulations, exporting the data to a structured CSV file, which a MATLAB script then parses using the \texttt{surf()} function and \texttt{exportgraphics} to produce high-resolution 3D surface plots.

These plots evaluate BB84 and B92 against channel Noise Strength ($p$) and eavesdropper Interception Rate ($\lambda$) simultaneously. Each $(p, \lambda)$ coordinate is the mean of $N_{\text{trials}}$ independent runs, ensuring the surfaces reflect reliable protocol behaviour rather than random simulation noise.

\section{Assumptions \& Limitations}
\label{sec:assumptionslimitations}

To scope the computational complexity of the project, the simulation relies on several necessary abstractions, detailed in \Cref{tab:limitations}.

\begin{table}[htbp]
\centering
\caption{Simulation Assumptions and Practical Implications.}
\label{tab:limitations}
\footnotesize
\begin{tabularx}{\linewidth}{@{}l X@{}}
\toprule
\textbf{Simulation Assumption} & \textbf{Practical Implication} \\
\midrule
\textbf{Shannon Capacity Limit} & Modelling classical error correction (Cascade/LDPC) is excluded. The Devetak-Winter bound assumes perfectly efficient reconciliation. In deployment, protocol inefficiencies would yield lower key rates than these theoretical bounds. \\
\addlinespace
\textbf{Ideal Hardware} & Qiskit's \texttt{AerSimulator} perfectly prepares single-photon states. Hardware imperfections (e.g., dark counts) are not modelled; the observed \gls{qber} is strictly a product of the depolarising channel and eavesdropping. \\
\addlinespace
\textbf{E91 Eavesdropper} & The \texttt{EveInterceptor} is constrained to \gls{pandm} protocols. \gls{e91} relies on distributed entanglement; intercepting one arm breaks the singlet state entirely. Simulating entanglement-swapping attacks is out of scope. \\
\addlinespace
\textbf{Perfect Classical Authenticity} & The classical channel (\gls{cc}) is assumed to be perfectly authenticated and immune to classical man-in-the-middle attacks during the sifting and post-processing phases. \\
\bottomrule
\end{tabularx}
\end{table}


%================================================================================
%                  Results 
%================================================================================
\chapter{Results \& Discussion}
%================================================================================

% overview
% remember to include reference to the other simualtion papers throughout to validate the results.

\section{Baseline Protocol Verification}
\label{sec:baseline}

To establish a control baseline, the protocols are evaluated over an ideal \gls{qchan}, absent of environmental noise or adversarial interference. As depicted in \Cref{fig:baseline}, the Mutual Information for all three protocols converges to the theoretical maximum, with the values of \gls{bb84} ($\sim$50), \gls{b92} ($\sim$25), and \gls{e91} ($\sim$22 bits/100 \glspl{qubit}) directly reflecting each protocol's sifting efficiency.


% mutual info bar chart
\begin{figure}[htbp]
    \centering
    \includegraphics[width=0.7\linewidth]{images.results v2/exp1/exp1_baseline_mutual_info.png}
    \caption{Baseline Mutual Information under Ideal Conditions.}
    \label{fig:baseline}
\end{figure}

\Cref{tab:validation} presents the theoretically expected values alongside the simulated means, where every simulated value lies within $\pm 2\%$ of its theoretical target, and tight error bars across all 50 trials confirm Monte Carlo stability. Perfect agreement across all seven metrics validates the simulation as a sound control environment prior to introducing interference models, aligning with the approach of recent \gls{qkd} simulation platforms \cite{devi2025}.

\begin{table}[htbp]
    \centering
    \caption{Ideal Channel Validation Results ($n=1000$, $N_{trials}=50$)}
    \label{tab:validation}
    \footnotesize
    \begin{tabular}{@{}l c c@{}}
        \toprule
        \textbf{Metric} & \textbf{Target} & \textbf{Simulated} \\
        \midrule
        BB84 Sifting Efficiency (\%) & $50.0$ & $50.0$ \\
        B92 Sifting Efficiency (\%)  & $25.0$ & $25.0$ \\
        E91 Sifting Efficiency (\%)  & $22.2$ & $22.2$ \\
        BB84 ideal-channel QBER (\%) & $0.0$  & $0.0$  \\
        B92 ideal-channel QBER (\%)  & $0.0$  & $0.0$  \\
        E91 ideal-channel QBER (\%)  & $0.0$  & $0.0$  \\
        E91 CHSH parameter \(|S|\)   & \(2\sqrt{2}\!\approx\!2.828\) & $2.828$ \\
        \bottomrule
    \end{tabular}
\end{table}

\newpage
\section{Resilience to Channel Noise}
\label{sec:resilience}

As modelled in \Cref{fig:qbernoise}, increasing Noise Strength directly elevates the \gls{qber} across all three protocols, with \gls{bb84} remaining below its $11\%$ security threshold up to $p \approx 0.22$. Notably, \gls{b92} and \gls{e91} exhibit near-identical \gls{qber} trajectories throughout, both breaching the \gls{b92} threshold of $6.5\%$ at $p \approx 0.06$, despite their different underlying security architectures. This \gls{bb84} behaviour is further validated by the results of Devi \textit{et al.}\ \cite{devi2025}, seen in \Cref{fig:deviqber}.


% \gls{qber} vs noise - increase
\begin{figure}[htbp]
    \centering
    \includegraphics[width=0.7\linewidth]{images.results v2/exp2/exp2_qber_vs_noise.png}
    \caption{QBER vs. Noise Strength, under the Depolarising Channel Model.}
    \label{fig:qbernoise}
\end{figure}

\begin{figure}[htbp]
    \centering
    \includegraphics[width=0.7\linewidth]{images.results v2//comparisons/bb84E91_noise_qber.png}
    \caption{Comparative QBER Protocol Performance Results for BB84 \& E91 Produced by Devi \textit{et al}.}
    \label{fig:deviqber}
\end{figure}

As shown in \Cref{fig:mutualnoise}, $I(A;B)$ degrades across all three protocols under increasing noise, with \gls{b92} and \gls{e91} converging toward similar residual capacities at high noise, whilst \gls{bb84} retains approximately twice their Mutual Information at $p = 0.30$.

% mutual info vs noise - degredation
\begin{figure}[htbp]
    \centering
    \includegraphics[width=0.7\linewidth]{images.results v2/exp2/exp2_mutual_info_vs_noise.png}
    \caption{Mutual Information vs. Noise Strength, under the Depolarising Channel Model.}
    \label{fig:mutualnoise}
\end{figure}

\newpage
Ultimately, the \gls{skr} consequences are shown in \Cref{fig:keyratenoise}. At $p = 5\%$, \gls{bb84} sustains $\sim\!34\ \text{bits/100 \glspl{qubit}}$, whilst \gls{e91} reduces to $\sim\!5$ and \gls{b92} collapses to $\sim\!4\ \text{bits/100 \glspl{qubit}}$. Compared against Kukutla \textit{et al.}\ \cite{kukutla2025}, \Cref{fig:deviskr} depicts the \gls{skr} for \gls{bb84} and \gls{e91}, where the disparity in magnitude reflects this framework's use of transmitted \glspl{qubit} rather than sifted \glspl{qubit} in the \gls{skr} denominator.

% \gls{skr}s vs noise - degredation
\begin{figure}[htbp]
    \centering
    \includegraphics[width=0.7\linewidth]{images.results v2/exp2/exp2_secure_key_rate_vs_noise.png}
    \caption{Secure Key Rate vs. Noise Strength, under the Depolarising Channel Model.}
    \label{fig:keyratenoise}
\end{figure}


%comparative results
\begin{figure}[htbp]
    \centering
    \includegraphics[width=0.7\linewidth]{images.results v2//comparisons/bb84E91_siftedkey_skr_noise.png}
    \caption{Comparative SKR Protocol Performance Results for BB84 \& E91 Produced by Kukutla \textit{et al}.}
    \label{fig:deviskr}
\end{figure}

The steeper \gls{b92} degradation stems from its two-state encoding, where depolarising noise introduces measurement errors that Bob cannot resolve, forcing a severe discard rate during \gls{sifting} \cite{B92}. Critically, despite sharing an almost identical \gls{qber} trajectory with \gls{b92}, \gls{e91} sustains a greater \gls{skr} above $p \approx 2.5\%$, with its Devetak-Winter threshold at $p \approx 0.120$ against \gls{b92}'s $p \approx 0.110$. This demonstrates that \gls{e91}'s \gls{eb} correlations provide a more favourable security penalty under the Devetak-Winter bound, even when the channel quality is equivalent.

\section{Eavesdropper Vulnerability}
\label{sec:everesults}

\Cref{fig:qberEve} evaluates \gls{bb84} and \gls{b92} under an active intercept-resend attack over an ideal \gls{qchan}, with horizontal threshold lines marking the Devetak-Winter bounds at which \gls{skr} drops to zero.

% \gls{qber} increase 2D plot
\begin{figure}[htbp]
    \centering
    \includegraphics[width=0.7\linewidth]{images.results v2/exp3/exp3_eve_vulnerability_qber.png}
    \caption{QBER vs. Intercept Rate, under an Ideal Channel Model.}
    \label{fig:qberEve}
\end{figure}



For \gls{bb84}, the $11\%$ security limit \cite{ShorPreskill2000} is breached at Eve $\approx 0.44$, confirming the established Devetak-Winter threshold. In contrast, \gls{b92} reaches its $6.5\%$ threshold at a mere Eve $\approx 0.13$, over three times more vulnerable to interception. Beyond their respective thresholds, both protocols must abort key generation. At full interception, \gls{b92} approaches $\sim\!50\%$ \gls{qber} against \gls{bb84}'s $\sim\!25\%$, consistent with the theoretical intercept-resend maxima for two-state and four-state encodings respectively \cite{ShorPreskill2000}.

The diverging gradient of the two curves above the \gls{b92} threshold is the key finding. \gls{b92}'s \gls{qber} ascent substantially steeper than \gls{bb84}'s at all Eve rates above $\approx 0.13$, demonstrating that the simplified two-state architecture compounds its lower threshold with accelerated error growth under sustained adversarial pressure.


\section{E91 Entanglement Mechanics}
\label{sec:e91results}

To verify the entanglement mechanics established in \Cref{subsec:CHSH}, this experiment models the \gls{e91} protocol under a depolarising channel. \Cref{fig:chshqber} presents a dual-axis evaluation adapted from the framework of Rasay \textit{et al.}\ \cite{rasay2026}, plotting the simulated S-parameter and \gls{qber} against increasing Noise Strength. The initial $|S| \approx 2.75$, slightly below the Tsirelson bound, reflects the both-side noise configuration applying depolarisation to both photons of the \gls{epr} pair simultaneously. The Kukutla \textit{et al.} comparison in \Cref{fig:kukchsh} \cite{kukutla2025} is broadly consistent, identifying the \gls{chsh} classical limit crossing at $\sim\!15\%$ noise against this platform's $P \approx 0.165$.


% \gls{chsh} \gls{qber} 2 axis plot
\begin{figure}[htbp]
    \centering
    \includegraphics[width=0.7\linewidth]{images.results v2/exp4/exp4_chsh_qber.png}
    \caption{Dual-Axis CHSH \& QBER vs. Noise Strength, under the Depolarising Channel Model.}
    \label{fig:chshqber}
\end{figure}

\begin{figure}[htbp]
    \centering
    \includegraphics[width=0.7\linewidth]{images.results v2//comparisons/e91CHSHvalidation_noise.png}
    \caption{Comparative Protocol Data for E91 Produced by Kukutla \textit{et al}.}
    \label{fig:kukchsh}
\end{figure}

As shown in \Cref{fig:chshqber}, the S-parameter degrades below the classical limit of $|S| = 2$ at $P \approx 0.165$, at which point $\text{QBER} \approx 17\%$ enables the inference of an approximate \gls{qber} threshold for this \gls{e91} configuration. The cryptographic consequence is quantified in \Cref{fig:e91keyrate}, which reveals that the Devetak-Winter bound forces \gls{skr} to zero at the earlier threshold of $P \approx 0.130$, preceding the full \gls{chsh} classical bound failure at $P \approx 0.165$ by a meaningful margin.

% \gls{skr} vs noise, overlaid with mutual info to highlight the separation between mutual info and usable key rate
\begin{figure}[htbp]
    \centering
    \includegraphics[width=0.7\linewidth]{images.results v2/exp4/exp4_secure_key_rate.png}
    \caption{E91 Secure Key Rate vs. Noise, under the Depolarising Channel Model, with Mutual Information Overlaid.}
    \label{fig:e91keyrate}
\end{figure}

As shown in \Cref{fig:e91keyrate}, \gls{skr} reaches zero at $P \approx 0.130$ whilst $I(A;B)$ remains substantially non-zero, confirming that the channel retains residual classical capacity well beyond the point at which secure key distillation has failed.

\newpage
\section{Multivariable System Dynamics}
\label{sec:multi}
In a real-world \gls{qchan}, noise and adversarial interception act simultaneously. To evaluate these multivariable dynamics, 3D surface    plots map Noise Strength and Eve's Intercept Rate against the \gls{qber}, Mutual Information, and \gls{skr} for \gls{bb84} and \gls{b92}. Threshold planes mark the volumetric boundary at which \gls{skr} drops to zero, with the Mutual Information surfaces establishing the theoretical security boundary that precedes it.


\subsection{BB84 Performance}
\label{subsec:bb84perf}

\subsubsection{QBER Surface}
\label{subsubsec:bb84qbersurface}

As previously described in \Cref{sec:protocolimplementation}, the security degradation of \gls{pandm} is evaluated by the \gls{qber} of a protocol. Hence, translating the methods of Rasay \textit{et al.} \cite{rasay2026} depicted in \Cref{fig:bb84rasayQBER} into three-dimensions, the \gls{qber} surface modelled in \Cref{fig:bb84qbersurface} depicts the operation of BB84, bounded by the well-established $11\%$ \gls{qber} threshold plane. 

% \gls{qber} surface plot
\begin{figure}[htbp]
    \centering
    \includegraphics[width=0.95\linewidth]{images.results v2/exp5/exp5_bb84_qber_v5.png}
    \caption{BB84 Multivariable QBER Surface.}
    \label{fig:bb84qbersurface}
\end{figure}

As shown in \Cref{tab:bb84QBERpoints}, even minimal baseline noise  accelerates the \gls{qber} trajectory, compressing the maximum tolerable intercept rate before the $11\%$ plane is breached. Nevertheless, \gls{bb84} retains a significant secure operational region beneath the threshold, validating key distribution in moderately noisy environments.

% \gls{qber} surface
% match data for the 6 points, 2 secure, 4 insecure
%%%%%%%%%%%%%%%%%%%%%%%%%%%%%%%%%%%%%%%%%%%%%%%%%%%%%%%%%%%%%%%%%%%%%%%%%%%%%%%%%%%%%%%%%%%%%%%%%%%%%%%%%%%%
\newpage
\begin{figure}[htbp]
    \centering
    \includegraphics[width=0.95\linewidth]{images.results v2//comparisons/bb84qberheatmap_rasay.png}
    \caption{BB84 QBER Heat-map Produced by Rasay \textit{et al.}}
    \label{fig:bb84rasayQBER}
\end{figure}


\begin{table}[htbp]
    \centering
    \caption{BB84 QBER Surface Samples ($n_{\gls{qubit}s}=500$, $N_{trials}=40$).}
    \label{tab:bb84QBERpoints}
    \footnotesize
    \begin{tabular}{@{}l c c c@{}}
        \toprule
        \textbf{Eve Rate (\%)} & \textbf{Noise Strength (\%)} & \textbf{QBER (\%)} & \textbf{Security}\\
        \midrule
        $7.50\%$ & $8.00\%$ & \textbf{$5.43\%$} & \textbf{Secure}\\
        $22.5\%$ & $4.00\%$ & \textbf{$7.26\%$} & \textbf{Secure}\\
        $6.00\%$ & $20.0\%$ & \textbf{$11.3\%$} & Insecure\\
        $30.0\%$ & $13.0\%$ & \textbf{$12.9\%$} & Insecure\\
        $39.0\%$ & $21.0\%$ & \textbf{$18.5\%$} & Insecure\\
        $46.5\%$ & $24.0\%$ & \textbf{$21.3\%$} & Insecure\\
        \bottomrule
    \end{tabular}
\end{table}

%%%%%%%%%%%%%%%%%%%%%%%%%%%%%%%%%%%%%%%%%%%%%%%%%%%%%%%%%%%%%%%%%%%%%%%%%%%%%%%%%%%%%%%%%%%%%%%%%%%%%%%%%%%%


\newpage
\subsubsection{Mutual Information Surface}
\label{subsubsec:bb84MIsurface}

The surface model in \Cref{fig:bb84MIsurface} plots $I(A;B)$ and $I(A;E)$ simultaneously, with the intersection ridge marking the security boundary $I_{AB} = I_{AE}$. As shown in \Cref{tab:bb84MIsurface}, $I(A;E)$ rises steeply with Eve's interception rate, with $I(A;B)$ also degrading under noise. This asymmetry causes the intersection ridge to move towards lower intercept rates as Noise Strength increases, narrowing the secure operating region under combined interference. Notably, $I(A;B)$ remains non-zero well beyond the security boundary, confirming that the channel retains residual classical capacity after \gls{skr} has already collapsed to zero.

% MI surface plot
\begin{figure}[htbp]
    \centering
    \includegraphics[width=0.95\linewidth]{images.results v2/exp5/exp5_bb84_mutual_info_v5.png}
    \caption{BB84 Multivariable Mutual Information Surfaces.}
    \label{fig:bb84MIsurface}
\end{figure}

% MI surface
\begin{table}[htbp]
    \centering
    \caption{BB84 Mutual Information Surface Samples ($n_{\gls{qubit}s}=500$, $N_{trials}=40$).}
    \label{tab:bb84MIsurface}
    \footnotesize
    \begin{tabular}{@{}l c c c@{}}
        \toprule
        \textbf{Eve Rate (\%)} & \textbf{Noise Strength (\%)} & \textbf{Mutual Information (bits per 100 \glspl{qubit})} & \textbf{Security}\\
        \midrule
        $4.5\%$ & $4.00\%$ & \textbf{$39.5$} & \textbf{Secure}  \\
        $24.0\%$ & $3.00\%$ & \textbf{$31.4$} & \textbf{Secure}  \\
        $3.00\%$ & $16.0\%$ & \textbf{$28.6$} & \textbf{Secure}  \\
        $10.5\%$ & $22.0\%$ & \textbf{$22.6$} & Insecure  \\
        $46.5\%$ & $20.0\%$ & \textbf{$15.2$} & Insecure  \\
        $49.5\%$ & $26.0\%$ & \textbf{$11.9$} & Insecure  \\
        \bottomrule
    \end{tabular}
\end{table}

\newpage
\subsubsection{Secure Key Rate Surface}
\label{subsubsec:bb84SKRsurface}

As shown in \Cref{fig:bb84SKRsurface}, \gls{skr} decays non-linearly from its peak at minimal noise and interception, collapsing to the flat zero plane across the majority of the operating surface. The data in \Cref{tab:bb84KRsurface} confirms that this collapse is accelerated by the interaction between the two interference sources convolving to drive \gls{skr} down to zero significantly faster than either source in isolation. This is consistent with the MI surface in \Cref{fig:bb84MIsurface}, where the intersection ridge migrated under combined interference, confirming that the Devetak-Winter bound absorbs the compounding penalty of both threats simultaneously.

% \gls{skr} surface plot
\begin{figure}[htbp]
    \centering
    \includegraphics[width=0.95\linewidth]{images.results v2/exp5/exp5_bb84_skr_v5.png}
    \caption{BB84 Multivariable Secure Key Rate Surface.}
    \label{fig:bb84SKRsurface}
\end{figure}

% \gls{skr} surface
\begin{table}[htbp]
    \centering
    \caption{BB84 Secure Key Rate Surface Samples ($n_{\gls{qubit}s}=500$, $N_{trials}=40$).}
    \label{tab:bb84KRsurface}
    \footnotesize
    \begin{tabular}{@{}l c c@{}}
        \toprule
        \textbf{Eve Rate (\%)} & \textbf{Noise Strength (\%)} & \textbf{\gls{skr} (bits per 100 \glspl{qubit})}\\
        \midrule
        $6.00\%$ & $4.00\%$ & \textbf{$27.8$} \\
        $4.50\%$ & $9.00\%$ & \textbf{$18.0$} \\
        $12.0\%$ & $13.0\%$ & \textbf{$6.12$} \\
        $28.5\%$ & $6.00\%$ & \textbf{$4.48$} \\
        $6.00\%$ & $19.0\%$ & \textbf{$2.22$} \\
        $36.0\%$ & $17.0\%$ & \textbf{$0.00$} \\    
        \bottomrule
    \end{tabular}
\end{table}

%%%%%%%%%%%%%%%%%%                           \gls{b92}                                %%%%%%%%%%%%%%%%%%%
\newpage
\subsection{B92 Performance}
\label{subsec:b92perf}

\subsubsection{QBER Surface}
\label{subsubsec:b92QBERsurface}

As shown in \Cref{fig:b92QBERsurface}, \gls{b92}'s stricter $6.5\%$ threshold is breached at substantially lower interference levels than \gls{bb84}, with the secure operational volume confined to a small corner of the parameter space. The surface ascends most steeply along the Eve Interception Rate axis, reflecting \gls{b92}'s two-state encoding mechanics, where inconclusive measurement events introduced by eavesdropping compound the sifting discard rate, producing a disproportionate \gls{qber} increase relative to noise. This asymmetry is consistent with the \gls{b92} \gls{qber} heat-map produced by Rasay \textit{et al.} \cite{rasay2026}, shown in \Cref{fig:b92rasayQBER}, which similarly identifies interception as the dominant degradation mechanism.

% \gls{qber} surface plot
\begin{figure}[htbp]
    \centering
    \includegraphics[width=0.95\linewidth]{images.results v2/exp6/exp6_b92_qber_v5.png}
    \caption{B92 Multivariable QBER Surface.}
    \label{fig:b92QBERsurface}
\end{figure}

Despite this fragility, \Cref{tab:b92QBERpoints} confirms that secure key distribution remains achievable under near-ideal conditions, despite the operating regime being substantially narrower than \gls{bb84}. The threshold plane is pierced at Eve rates as low as $\sim\!10\%$ even with minimal channel noise, demonstrating that \gls{b92}'s simplified architecture carries a significant practical cost in adversarial environments.

% \gls{qber} surface
% match data for the 6 points, 2 secure, 4 insecure, chnage to ten samples used from rasays results to validate againt in a new column 
%%%%%%%%%%%%%%%%%%%%%%%%%%%%%%%%%%%%%%%%%%%%%%%%%%%%%%%%%%%%%%%%%%%%%%%%%%%%%%%%%%%%%%%%%%%%%%%%%%%%%%%%%%%%
\newpage
\begin{figure}[htbp]
    \centering
    \includegraphics[width=0.7\linewidth]{images.results v2//comparisons/b92qberheatmap_rasay.png}
    \caption{B92 QBER Heat-map Produced by Rasay \textit{et al.}}
    \label{fig:b92rasayQBER}
\end{figure}


\begin{table}[htbp]
    \centering
    \caption{B92 QBER Surface Samples ($n_{\gls{qubit}s}=500$, $N_{trials}=40$).}
    \label{tab:b92QBERpoints}
    \footnotesize
    \begin{tabular}{@{}l c c c@{}}
        \toprule
        \textbf{Eve Rate (\%)} & \textbf{Noise Strength (\%)} & \textbf{QBER} & \textbf{Security}\\
        \midrule
        $2.25\%$ & $2.00\%$ & \textbf{$3.32\%$} & \textbf{Secure}  \\
        $9.00\%$ & $1.50\%$ & \textbf{$5.37\%$} & \textbf{Secure}  \\
        $1.50\%$ & $5.50\%$ & \textbf{$5.60\%$} & Insecure  \\
        $9.75\%$ & $3.50\%$ & \textbf{$7.84\%$} & Insecure  \\
        $6.75\%$ & $6.50\%$ & \textbf{$9.14\%$} & Insecure  \\
        $3.75\%$ & $11.0\%$ & \textbf{$10.8\%$} & Insecure  \\
        \bottomrule
    \end{tabular}
\end{table}
%%%%%%%%%%%%%%%%%%%%%%%%%%%%%%%%%%%%%%%%%%%%%%%%%%%%%%%%%%%%%%%%%%%%%%%%%%%%%%%%%%%%%%%%%%%%%%%%%%%%%%%%%%%%


\newpage
\subsubsection{Mutual Information Surface}
\label{subsubsec:b92MIsurface}

The \gls{b92} Mutual Information surfaces in \Cref{fig:b92MIsurface} reveal a markedly more compressed $I(A;B)$ ceiling relative to \gls{bb84}, consistent with \gls{b92}'s lower sifting efficiency reducing the volume of secure key material available. The intersection where $I(A;E)$ overtakes $I(A;B)$ occurs at substantially lower Eve rates, confirming that the security boundary is reached under far less adversarial pressure than \gls{bb84}.

% MI surface plot
\begin{figure}[htbp]
    \centering
    \includegraphics[width=0.95\linewidth]{images.results v2/exp6/exp6_b92_mutual_info_v5.png}
    \caption{B92 Multivariable Mutual Information Surfaces.}
    \label{fig:b92MIsurface}
\end{figure}

Notably, $I(A;E)$ exhibits a dome-like profile, rising steeply with interception rate before plateauing at high noise and Eve rate combinations. As shown in \Cref{tab:b92MIsurface}, $I(A;B)$ remains non-zero throughout the insecure region, mirroring the behaviour observed in \gls{bb84} and confirming that classical channel capacity persists well beyond the point at which secure key distillation has failed.

% MI surface
\begin{table}[htbp]
    \centering
    \caption{B92 Mutual Information Surface Samples ($n_{\gls{qubit}s}=500$, $N_{trials}=40$).}
    \label{tab:b92MIsurface}
    \footnotesize
    \begin{tabular}{@{}l c c c@{}}
        \toprule
        \textbf{Eve Rate (\%)} & \textbf{Noise Strength (\%)} & \textbf{Mutual Information (bits per 100 \glspl{qubit})} & \textbf{Security}\\
        \midrule
        $1.50\%$ & $1.00\%$ & \textbf{$22.1$} & \textbf{Secure}  \\
        $2.25\%$ & $3.00\%$ & \textbf{$20.3$} & \textbf{Secure}  \\
        $16.5\%$ & $2.00\%$ & \textbf{$13.9$} & Insecure  \\
        $13.5\%$ & $7.00\%$ & \textbf{$12.2$} & Insecure  \\
        $10.5\%$ & $12.5\%$ & \textbf{$11.9$} & Insecure  \\ 
        $24.0\%$ & $9.50\%$ & \textbf{$8.60$} & Insecure  \\ 
        \bottomrule
    \end{tabular}
\end{table}


\newpage
\subsubsection{Secure Key Rate Surface}
\label{subsubsec:b92SKRsurface}

As shown in \Cref{fig:b92SKRsurface}, \gls{skr} collapses to zero within an extremely narrow viable region, with the zero plane dominating the surface far more aggressively than \gls{bb84}. The peak \gls{skr} of $\sim\!25$ bits/100 transmitted qubits is approximately half that of \gls{bb84}, and the data in \Cref{tab:b92SKRsurface} confirms that even modest noise levels drive \gls{skr} to zero irrespective of Eve rate. Taken together, the three \gls{b92} surfaces empirically demonstrate that the protocol requires near-ideal channel conditions to sustain secure key generation, rendering it impractical for most real-world deployment environments.

% \gls{skr} surface plot
\begin{figure}[htbp]
    \centering
    \includegraphics[width=0.95\linewidth]{images.results v2/exp6/exp6_b92_skr_v5.png}
    \caption{B92 Multivariable Secure Key Rate Surface.}
    \label{fig:b92SKRsurface}
\end{figure}

% \gls{skr} surface
\begin{table}[htbp]
    \centering
    \caption{B92 Secure Key Rate Surface Samples ($n_{\gls{qubit}s}=500$, $N_{trials}=40$).}
    \label{tab:b92SKRsurface}
    \footnotesize
    \begin{tabular}{@{}l c c@{}}
        \toprule
        \textbf{Eve Rate (\%)} & \textbf{Noise Strength (\%)} & \textbf{\gls{skr} (bits per 100 \glspl{qubit})}\\
        \midrule
        $1.50\%$ & $0.50\%$ & \textbf{$17.9$} \\
        $1.50\%$ & $2.00\%$ & \textbf{$12.3$} \\
        $5.25\%$ & $1.50\%$ & \textbf{$7.25$} \\
        $5.25\%$ & $3.00\%$ & \textbf{$3.25$} \\
        $9.75\%$ & $1.00\%$ & \textbf{$1.93$} \\
        $21.0\%$ & $7.00\%$ & \textbf{$0.00$} \\
        \bottomrule
    \end{tabular}
\end{table}

It should be noted that in practical deployment, distinguishing between channel noise and adversarial interception is notoriously difficult, even with an established operational noise regime known to both parties prior to transmission. Hence, the worst-case attribution scheme applied throughout (\Cref{subsec:DWsec}) represents the conservative but necessary bound for all \gls{skr} calculations presented here. Despite this, the graphs provide a useful visual aid for understanding the multivariable system dynamics of \gls{bb84} and \gls{b92}, advancing the work of Rasay \textit{et al.} \cite{rasay2026} with novel representations of the multivariable system dynamics of \gls{bb84} and \gls{b92}.


%================================================================================
%                  Conclusion 
%================================================================================
\chapter{Conclusion \& Future Work}
%================================================================================

This technical report detailed the design, implementation, and evaluation of a \gls{qkd} simulation platform. Using a modular \gls{oop} architecture, the project successfully modelled three \gls{dv}-\gls{qkd} protocols under realistic channel conditions, translating abstract quantum mechanics into quantifiable engineering metrics. This chapter evaluates the project's achievements against its initial objectives, synthesises key findings, and outlines the broader societal context and directions for future research.

\section{Achievement Against Objectives}
\label{sec:achievements}

The primary aim of developing a configurable, software-based \gls{qkd} simulation platform to benchmark protocol performance under realistic constraints was successfully achieved. The modular architecture directly fulfilled the implementation and modelling objectives (\textbf{O2, O3}), successfully simulating state preparation, depolarising channel noise, and active intercept-resend attacks across the \gls{bb84}, \gls{b92}, and \gls{e91} protocols. 

The integration of a rigorous analytical framework (\textbf{O1}) and a comprehensive post-processing pipeline enabled scalable, multivariable batch testing. This facilitated accurate error estimation and the calculation of Devetak-Winter secure key rates (\textbf{O4}), allowing for the systematic validation of protocol vulnerabilities against established theoretical bounds.


\section{Key Findings}
\label{sec:key}

\Cref{tab:key_findings} details the critical findings yielded upon completion of this project.

\begin{table}[htbp]
    \centering
    \caption{Critical Findings from the Multivariable Protocol Analysis.}
    \label{tab:key_findings}
    \footnotesize
    \renewcommand{\arraystretch}{1.4} % Improves readability for dense text
    \begin{tabular}{@{}p{0.25\linewidth} p{0.7\linewidth}@{}}
        \toprule
        \textbf{Analytical Focus} & \textbf{Key Result \& Observation} \\
        \midrule
        \textbf{Architecture Validation} & Under ideal channel conditions, the simulation converged to the theoretical Mutual Information baseline of \( 1.0\ \text{bit/sifted qubit} \) with a \( 0.0\% \) \gls{qber}, validating the core cryptographic logic. \\
        
        \textbf{\gls{bb84} vs. \gls{b92} Resilience} & The four-state structure of \gls{bb84} demonstrated significant resilience to depolarising noise, maintaining a \gls{skr} of \( \sim34\ \text{bits/100 \gls{qubit}s} \) at \( 5\% \) noise. \gls{b92}'s two-state structure caused a severe sifting discard rate, collapsing the \gls{skr} to \( \sim4\ \text{bits/100 \gls{qubit}s} \). \\
        
        \textbf{Entanglement Degradation} & Successfully mapped the threshold of \gls{e91} \gls{entanglement} collapse. At a noise strength of \( P \approx 0.165 \), the S-parameter decayed from the quantum Tsirelson bound of \( 2\sqrt{2} \) to below the classical limit of \( 2 \), resulting in key distribution failure at \( P \approx 0.130 \). \\
        
        \textbf{Deployment Viability} & 3D surface analysis proved that while \gls{bb84} maintains a functional secure operational volume under simultaneous noise and interception, \gls{b92} rapidly breaches its \( 6.5\% \) \gls{qber} threshold, rendering its deployment largely impractical. \\
        \bottomrule
    \end{tabular}
\end{table}

\section{Evaluation of Methodology \& Limitations}
\label{sec:eval}
While Qiskit and NumPy facilitate theoretical \gls{qkd} analysis, they abstract physical device imperfections such as detector inefficiency and fibre attenuation. Furthermore, the simulation currently excludes classical error correction algorithms like Cascade \cite{Cascade1994} or LDPC. By assuming perfectly efficient reconciliation at the Shannon limit, the multivariable SKR surfaces represent an absolute theoretical upper bound. In a physical deployment, the reconciliation efficiency factor (\(f > 1\)) would compress these SKR volumes closer to the origin, narrowing the operational viability of both \gls{bb84} and \gls{b92}.

A secondary limitation exists within the active eavesdropping implementation, which models the intercept-resend strategy. For B92 specifically, an Unambiguous State Discrimination (USD) attack would exploit its two-state encoding more effectively during the physical transmission phase. However, because the system's SKR calculations evaluate the Devetak-Winter bound for optimal attacks, the classical post-processing pipeline applies the maximum theoretical penalty regardless of the simulated attack model. This ensures that the generated performance thresholds and secure key estimations remain mathematically conservative, even if the simulated intercept-resend attack underperforms a true USD attack.

Although Monte Carlo methods and the Central Limit Theorem mitigate statistical uncertainty across the parameter sweeps, the derived metrics must be interpreted as theoretical limits rather than physical benchmarks. Overall, the project reflects deliberate engineering compromises, abstracting hardware complexity to ensure the successful creation of a functional protocol evaluation platform.

\section{Broader Context, Ethics \& EDI}
\label{sec:context}

As quantum threats compromise classical cryptography, bridging quantum theory with practical network engineering is critical. Ethically, quantum-secure infrastructure must protect all sectors equally. From an \gls{edi} perspective, developing the simulator using open-source libraries and modular documentation lowers the barrier to entry, ensuring global accessibility for diverse researchers.

\section{Future Work}
\label{sec:future}

The following developments would advance the platform's deployment 
realism, driving the project further towards digital-twin systems engineering capability.

\begin{table}[htbp]
    \centering
    \caption{Proposed Directions for Future Work}
    \label{tab:future_work}
    \renewcommand{\arraystretch}{1.4} % Adds vertical padding for readability
    \footnotesize
    \begin{tabular}{@{}p{0.3\linewidth} p{0.65\linewidth}@{}}
        \toprule
        \textbf{Proposed Development} & \textbf{Technical Objective \& Justification} \\
        \midrule
        \textbf{Implementation of Eve for \gls{e91}} & Unifying the benchmark by developing an optimal attack strategy against \gls{e91} to complete the comprehensive analysis of the three protocols. \\
        
        \textbf{Advanced \gls{qkd} Protocols} & Integrating alternative protocols described by Kumar \textit{et al.} \cite{Kumar2025} to expand the simulator suite to modern protocols. \\
        
        \textbf{Hardware Noise Integration} & Coupling the simulator with IBM Quantum real hardware noise models to evolve the channel interference module from theoretical approximations to device-specific error characteristics. \\
        
        \textbf{Expanded Network Topology} & Advancing the architecture from simple \gls{p2p} links to multi-node quantum network topologies to enable the simulation of quantum repeaters and trusted-node routing. \\
        \bottomrule
    \end{tabular}
\end{table}

\newpage
\section{Generative AI Usage}
\label{sec:ai_declaration}

Generative AI assisted with conceptual learning, literature scoping, and grammatical refinement. All software architecture and analysis remain independent work. \Cref{tab:ai_usage} details specific AI usage and oversight.

\begin{table}[htbp]
    \centering
    \caption{Summary of Generative AI Usage and Compliance Oversight}
    \label{tab:ai_usage}
    \footnotesize
    \begin{tabularx}{\linewidth}{@{}l l X X@{}}
        \toprule
        \textbf{Project Phase} & \textbf{Tool(s) Used} & \textbf{Nature of Use} & \textbf{Oversight \& Compliance} \\
        \midrule
        \textbf{Conceptual Learning} & Claude & Acted as an interactive study tool to break down abstract quantum mechanics, Bell inequalities, and classical post-processing equations. & Explanations were manually verified against foundational textbooks (Nielsen \cite{QKDbible}, Ekert \cite{ekertQI}) to ensure mathematical rigour. \\
        \addlinespace
        \textbf{Literature Survey} & Perplexity, Google Scholar Labs & Used to conduct initial literature scoping, identify relevant academic papers on \gls{qkd} simulation. & Primary sources were read manually. All literature analysis, and discussion of findings within the report were conducted independently. \\
        \addlinespace
        \textbf{Project Planning} & Perplexity & Provided insights into how to plan and segment stages of the project to improve time management. & These plans informed setting SMART objectives throughout the course of the project. \\
        \addlinespace
        \textbf{Software Development} & Claude & Used to clarify Qiskit syntax, understand backend measurement mechanics, and assist with debugging Python \& MATLAB errors. & All core \gls{oop} architecture, protocol logic, and experimental designs were a product of independent work. AI was limited to syntax learning, not architectural generation. \\
        \addlinespace
        \textbf{Drafting \& Formatting} & Perplexity & Aided with identifying grammatical errors, guidance on LaTeX code preamble formatting, and to check academic tone. & The original meaning and technical arguments of the project are also a product of independent work. \\
        \bottomrule
    \end{tabularx}
\end{table}

While capable of summarising classical programming concepts, LLMs lacked the mathematical depth required for quantum simulation and frequently omitted critical noise models in literature reviews. Manual verification ensured that theoretical accuracy was achieved.




% remove before compile
%\chapter{ }
% ─────────────────────────────────────────────────────────────────────────────
% REFERENCES
% ─────────────────────────────────────────────────────────────────────────────
% IEEE references: 8pt, 9pt line spacing, bracketed Arabic numerals [1], [2]
\clearpage
{\footnotesize
\bibliographystyle{IEEEtran}
\bibliography{bib}
}

%55555555555555555555555555555555555555555555555555555
\appendix
\chapter{Source Code}
% Full source code listings or link to repository.
\label{app:source_code}

The complete Python source code for the \gls{qkd} simulation platform developed for this project is hosted in a public GitHub repository. This repo contains the protocol implementations (BB84, B92, E91), the experiment configuration system, and the MATLAB visualisation scripts. 

The repository can be accessed via the following link:

\vspace{0.5cm}
\begin{center}
    \url{https://github.com/omdillon/QKD}
\end{center}
\vspace{0.5cm}

Full instructions for configuring the Python environment, installing the necessary dependencies (including Qiskit and NumPy), and executing the simulation modes are documented within the repository's \texttt{README.md} file.

\end{document}